\documentclass[11pt,a4paper]{article}

% --- Packages ---
\usepackage[margin=2.5cm]{geometry}
\setlength{\headheight}{14pt}
\usepackage{hyperref}
\usepackage{titlesec}
\usepackage{enumitem}
\usepackage{booktabs}
\usepackage{longtable}
\usepackage{xcolor}
\usepackage{fancyhdr}
\usepackage{tikz}
\usetikzlibrary{shapes.geometric, arrows.meta, positioning, fit, backgrounds}
\usepackage{listings}
\usepackage{parskip}
\usepackage{microtype}
\usepackage{amsmath}
\usepackage{amssymb}
\usepackage{mathtools}

% --- Style ---
\hypersetup{colorlinks=true, linkcolor=blue!60!black, urlcolor=blue!60!black}
\titleformat{\section}{\large\bfseries}{}{0em}{}[\titlerule]
\titleformat{\subsection}{\normalsize\bfseries}{}{0em}{}
\titleformat{\subsubsection}{\normalsize\itshape}{}{0em}{}
\pagestyle{fancy}
\fancyhf{}
\rhead{\textit{Fleet Management --- SRS}}
%% --- CHANGE RF-001: version bump to 0.4 (running header) ---
\lhead{\textit{v0.4 $\cdot$ \today}}
\cfoot{\thepage}

\definecolor{todo}{RGB}{180,0,0}
\newcommand{\TODO}[1]{\textcolor{todo}{\textbf{[TO DO: #1]}}}

\lstset{
  basicstyle=\ttfamily\small,
  breaklines=true,
  frame=single,
  backgroundcolor=\color{gray!8},
  rulecolor=\color{gray!40}
}

\begin{document}

% --- Title page ---
\begin{titlepage}
  \centering
  \vspace*{3cm}
  {\Huge\bfseries Fleet Management\par}
  \vspace{0.5cm}
  {\Large Software Requirements Specification\par}
  \vspace{0.3cm}
  {\large Multi-Component Stochastic Degradation MILP\par}
  \vspace{1cm}
  %% --- CHANGE RF-002: version bump to 0.4 (title page) ---
  {\large Version 0.4 $\cdot$ \today\par}
  \vfill
  {\small Prepared with spec-creator $\cdot$ Intended for Claude Code implementation}
\end{titlepage}

\tableofcontents
\newpage

% ================================================================
\section{Project Overview \& Goals}
% ================================================================

This project implements a \textbf{Mixed-Integer (Non-)Linear Program} for train fleet
scheduling with multi-component stochastic degradation, solved with Gurobi.

The key generalisation over prior work is that each component $\ell$ of each train $i$
may have its own degradation model, independently chosen from:
\begin{itemize}
  \item \textbf{Gaussian} (Gaussian random walk; tracks mean and variance independently)
  \item \textbf{Inverse Gaussian (IG)} (tracks mean only; shape parameter recovered analytically)
  \item \textbf{Wiener} (Brownian motion with drift; tracks mean and variance, but variance
        is a deterministic function of time with no repair effect)
  \item \textbf{Gamma} (Gamma random walk; tracks accumulated shape parameter)
  %% --- CHANGE RF-003: add rainflow model bullet ---
  \item \textbf{Rainflow} (distribution-free; the per-mission damage-increment law
        is measured empirically by rainflow counting of load histories and enters
        the model only through its mean, variance, and---for the Bernstein
        certificate---support upper bound; tracks mean and variance; reliability
        certified by a Cantelli or Bernstein tail bound)
\end{itemize}

%% --- CHANGE DOM-003: clarify Gaussian vs Wiener distinction ---
\emph{Gaussian vs.\ Wiener distinction.}
Both the Gaussian and Wiener models use Gaussian-distributed degradation increments.
The distinction is that the Gaussian model tracks both the mean~$\mu$ and the
variance~$v$ of each component as independent decision variables, allowing imperfect
maintenance to act on both quantities. The Wiener model, by contrast, treats variance
as a deterministic function of time ($v_k = v_0 + k\sigma^2$), with no repair effect
on variance. This makes the Wiener model simpler but restricts it to components where
periodic full replacement is feasible (see Section~\ref{sec:dyn_wiener}). When in
doubt, use the Gaussian model; the Wiener model is a specialisation for components
with stationary, route-independent uncertainty.

%% --- CHANGE RF-004: rainflow vs parametric models contrast remark ---
\emph{Rainflow vs.\ the parametric models.}
The rainflow model makes \textbf{no distributional assumption} on the damage
increments: the per-mission increment law is measured (rainflow counting of
recorded load histories) and is summarised only by its first two moments and,
optionally, its support upper bound. Reliability
$\Pr(D > \tau) \leq \varepsilon$ is certified by a distribution-free tail bound
rather than by an exact parametric CDF, trading tightness for robustness: the
certificate holds for \emph{every} increment law consistent with the measured
summaries. Two bounds are offered via the \texttt{tail\_bound} input
(Section~\ref{sec:dyn_rainflow}): \textbf{Cantelli}, which needs only the mean
and variance and is the always-valid coarse option, and \textbf{Bernstein},
which additionally exploits the mutual independence of the daily increments and
their bounded support, and is tighter deep in the tail at the price of requiring
support inputs.

Each component also has an independently specified maintenance policy
($\mathrm{ARA}_1$ or $\mathrm{ARD}_1$, subject to model compatibility---see below)
and may be subject to full replacement.

%% --- CHANGE DOM-008: ARD1/ARA1 definitions ---
\textbf{Maintenance policy definitions.}
The two imperfect-repair models used in this specification are:
\begin{itemize}
  \item $\mathrm{ARD}_1$ (\emph{Arithmetic Reduction of Degradation}, order 1):
        after repair, the degradation level is reduced by a fraction~$\rho$ of its
        current value, i.e.\ $D^+ = (1-\rho)\,D^-$.
        This is equivalent to removing a fixed proportion of accumulated damage
        (Doyen \& Gaudoin, 2004; Kijima, 1989).
  \item $\mathrm{ARA}_1$ (\emph{Arithmetic Reduction of Age}, order 1):
        after repair, the degradation level is reduced by a fraction~$\rho$ of the
        degradation accumulated since the last maintenance epoch, i.e.\
        $D^+ = D^- - \rho\,(D^- - D^{\mathrm{last}})$,
        where $D^{\mathrm{last}}$ is the degradation immediately after the most recent
        maintenance action (Mercier \& Castro, 2019; Kijima, 1989).
\end{itemize}

%% --- CHANGE DEV-002: prominently flag Wiener replacement requirement ---
\textbf{Wiener model restriction.}
Because imperfect maintenance cannot reduce variance in the Wiener model, every
Wiener-modelled component \textbf{must} undergo at least one full replacement
($x^r_{i\ell k}=1$) per horizon~$H$. Without this, the loop constraint on variance
is infeasible. This restriction should be verified at input validation time
(Section~3.7).

\textbf{Non-goals.}
\begin{itemize}
  \item Routing decisions are not optimised (routes are an input).
  \item Continuous-time formulations are out of scope; all dynamics are discrete-time.
  %% --- CHANGE RF-005: non-goals carve-out for rainflow linearisation ---
  %% --- CHANGE RF-042 (panel): bullet reworded; linearisation is no longer deferred anywhere ---
  \item Tightening of the linear reliability surrogates beyond their single-tangent
        forms is deferred. Every model ships with fully specified constraint forms
        in this document: Gaussian, Wiener, and rainflow have both an exact
        (nonconvex quadratic) and a linear inner-approximation form
        (Sections~\ref{sec:dyn_gaussian} and~\ref{sec:dyn_rainflow}); IG and Gamma
        are linear by construction. The deferred item is the multi-tangent PWL
        refinement listed in the Open Questions.
\end{itemize}

% ================================================================
\section{Index Sets and Notation}
% ================================================================

%% --- CHANGE MATH-001: j in {0,1,...,M} ---
\begin{longtable}{@{}lp{10cm}@{}}
\toprule
Symbol & Meaning \\
\midrule
$i \in \{1,\ldots,F\}$ & Train index \\
$j \in \{0,1,\ldots,M\}$ & Activity index ($j=0$: maintenance day; $j \in \{1,\ldots,M\}$: missions) \\
$k \in \{1,\ldots,2H\}$ & Time step (day) index \\
$\ell \in \{1,\ldots,L\}$ & Component index \\
\bottomrule
\end{longtable}

% ================================================================
\section{Input Specification}
% ================================================================

\subsection{Scalar / Fleet-Level Parameters}
\label{sec:input_scalars}

\begin{longtable}{@{}llp{8cm}@{}}
\toprule
Parameter & Type & Description \\
\midrule
$F$ & \texttt{int} & Number of trains \\
$H$ & \texttt{int} or \texttt{[int, int]} & Half-horizon. Scalar: single solve spanning $2H$
  steps. Two-element list $[H_{\min}, H_{\max}]$: loop over
  $H \in \{H_{\min},\ldots,H_{\max}\}$, producing one solve per value \\
$M$ & \texttt{int} & Number of missions (must satisfy $F > M$) \\
$L$ & \texttt{int} & Number of components per train \\
$\tau$ & \texttt{float} & Degradation upper bound (must be positive) \\
$\varepsilon$ & \texttt{float} & Reliability threshold, $\varepsilon \in (0, 0.01]$.
  The upper bound $\varepsilon \leq 0.01$ prevents $\Phi^{-1}(1-\varepsilon)$ from
  approaching zero, which would cause $V_{\max} \to \infty$ and render the
  big-M constraints numerically ineffective \\
$C_M$ & \texttt{float} & Maintenance cost coefficient \\
$C_P$ & \texttt{float} & Loop-constraint penalty coefficient \\
$C_D$ & \texttt{float} & Damage regularisation coefficient \\
\texttt{penalty\_type} & \texttt{str}, optional & Damage regularisation type:
  \texttt{"inf\_norm"} (default) or \texttt{"quadratic"}.
  Controls the form of the $C_D$ penalty term (see Section~\ref{sec:objective}) \\
%% --- CHANGE RF-006: formulation token socp->exact; corrected SOC mislabel ---
\texttt{formulation} & \texttt{str}, optional & Problem formulation:
  \texttt{"exact"} (default) or \texttt{"lp"}.
  \texttt{"exact"}: Gaussian, Wiener, and rainflow reliability constraints are
  imposed in their exact squared form, a \emph{nonconvex} quadratic constraint
  handled by Gurobi's nonconvex quadratic support (\texttt{NonConvex=2} where
  not the default); see Section~\ref{sec:dyn_gaussian}.
  \texttt{"lp"}: all constraints are linear; Gaussian/Wiener reliability
  constraints use the inner approximation
  $(\Phi^{-1}(1-\varepsilon))^2 v_{i\ell k} + 2\tau\,\mu_{i\ell k} \leq \tau^2$
  (conservative; approximation error $= \mu_{i\ell k}^2$, small when $\mu_{i\ell k} \ll \tau$),
  and rainflow components use the analogous tangent surrogates of
  Section~\ref{sec:dyn_rainflow}.
  \texttt{"quadratic"} penalty type is incompatible with \texttt{"lp"}.
  The legacy value \texttt{"socp"} is accepted as a deprecated alias for
  \texttt{"exact"}: the parser maps it to \texttt{"exact"} and emits a
  \texttt{UserWarning} \\
\texttt{n\_workers} & \texttt{int}, optional & Number of parallel Gurobi solves.
  \texttt{1} = sequential (default). Each worker runs an independent Gurobi environment.
  Only meaningful when $H$ is an interval \\
\texttt{warm\_start} & \texttt{bool}, optional & If \texttt{True}, pass the solution for
  $H_k$ as a heuristic hint (\texttt{VarHintVal}) for $H_{k+1}$. No feasibility guarantee.
  Default: \texttt{False}. Requires \texttt{n\_workers=1} \\
\bottomrule
\end{longtable}

\subsection{Component Model Assignment}

\begin{longtable}{@{}llp{8cm}@{}}
\toprule
Parameter & Shape & Description \\
\midrule
%% --- CHANGE RF-008: add rainflow to model strings ---
\texttt{model} & $F \times L$ & Degradation model per train-component:
  one of \texttt{"gaussian"}, \texttt{"inverse\_gaussian"},
  \texttt{"wiener"}, \texttt{"gamma"}, \texttt{"rainflow"} \\
%% --- CHANGE RF-007: new tail_bound input ---
%% --- CHANGE RF-044 (panel): tail_bound moved here from the scalar-parameter table (conditionally required, per-component) ---
\texttt{tail\_bound} & \texttt{str} or $F \times L$ list of \texttt{str} &
  Distribution-free tail bound used in the reliability constraint of rainflow
  components: \texttt{"cantelli"} (default) or \texttt{"bernstein"}, either a
  single string applied to all rainflow components or an $F \times L$
  per-component list. Optional; read only when at least one component is
  \texttt{"rainflow"}. Entries for non-rainflow components are ignored with a
  \texttt{UserWarning}. See Section~\ref{sec:dyn_rainflow} \\
\bottomrule
\end{longtable}

\subsection{Initial Conditions}

\begin{longtable}{@{}llp{8cm}@{}}
\toprule
Parameter & Shape & Description \\
\midrule
$\mu_0$ & $F \times L$ & Initial mean degradation (all models); must satisfy $\mu_0 < \tau$ \\
%% --- CHANGE RF-009: v_0 also applies to rainflow ---
$v_0$ & $F \times L$ & Initial variance (Gaussian, Wiener, and Rainflow); must satisfy
  $\mu_0 \geq 3\sqrt{v_0}$ for Gaussian \\
\bottomrule
\end{longtable}

\subsection{Mission- and Time-Dependent Degradation Parameters}
\label{sec:input_increments}

These tensors provide the degradation \emph{increment} parameters when train $i$
performs mission $j$ on day $k$ for component $\ell$.
Days $k > H$ reuse the values at $k - H$ (periodic extension).

%% --- CHANGE CS-007: specify eager broadcasting strategy ---
\emph{Broadcasting strategy.}
The solver shall use \textbf{eager broadcasting}: upon loading, all input parameters
of shape $(F, M, L)$ or $(F, M, L, H)$ are immediately expanded to shape
$(F, M, L, 2H)$ using periodic wrapping at index $k \bmod H$ for $k \geq H$.
This expansion happens once, before constraint construction, so that all
constraint-building code can assume full $(F, M, L, 2H)$ arrays.

\begin{longtable}{@{}lllp{5cm}@{}}
\toprule
Parameter & Shape & Models & Description \\
\midrule
%% --- CHANGE RF-010: mu_inc/v_inc also for rainflow; new b_inc row ---
$\mu^{\mathrm{inc}}$ & $F \times M \times L \times H$ & Gaussian, IG, Wiener, Rainflow
  & Mean increment per step \\
$v^{\mathrm{inc}}$ & $F \times M \times L \times H$ & Gaussian, Rainflow
  & Variance increment per step \\
$\alpha^{\mathrm{inc}}$ & $F \times M \times L \times H$ & Gamma only
  & Shape increment per step \\
$b^{\mathrm{inc}}$ & $F \times M \times L \times H$ & Rainflow (Bernstein only)
  & Support upper bound of the empirical per-mission damage increment, from
    rainflow counting; YAML key \texttt{b\_inc} \\
\bottomrule
\end{longtable}

\subsection{Component-Level Constants}

\begin{longtable}{@{}llp{8cm}@{}}
\toprule
Parameter & Shape & Description \\
\midrule
$\eta$ & $F \times L$ & $\mu/\lambda$ ratio (IG only; must be positive).
  The constant-$\eta$ assumption (proportional dispersion) is well-supported
  for degradation mechanisms exhibiting self-similar growth: in fatigue crack
  propagation (Paris' law) the coefficient of variation of increments is
  approximately constant, and many wear/corrosion processes share this property.
  This assumption enables the FSD reduction to a mean-only loop condition
  (Section~\ref{sec:dyn_ig}) and is preserved under ARD1 repair since both
  $\mu$ and $\lambda = \mu/\eta$ scale by the same factor $(1-\rho)$ \\
$\sigma$ & $F \times L$ & Diffusion constant (Wiener only; must be positive) \\
$\beta$ & $F \times L$ & Scale parameter (Gamma only; must be positive) \\
%% --- CHANGE DOM-026: user-defined variance upper bound ---
%% --- CHANGE RF-011: V_max^user applicability extended to rainflow ---
$V_{\max}^{\mathrm{user}}$ & $F \times L$, optional & User-defined variance upper bound
  (Gaussian, Wiener, and Rainflow components only). When specified, the constraint
  $v_{i\ell k} \leq V_{\max,i\ell}^{\mathrm{user}}$ is added for all $k$.
  Default: not applied (variance is bounded only by the reliability constraint) \\
\bottomrule
\end{longtable}

\subsection{Maintenance and Replacement Specification}

\begin{longtable}{@{}llp{8cm}@{}}
\toprule
Parameter & Shape & Description \\
\midrule
\texttt{maintenance\_type} & $F \times L$ & Per train-component maintenance type:
  \texttt{"ARA1"} or \texttt{"ARD1"} \\
$\rho$ & $F \times L$ & Repair efficiency $\in (0,1]$ \\
$C_R$ & $F \times L$ & Repair cost (ARA/ARD maintenance) \\
$C_{\mathrm{rep}}$ & $F \times L$ & Replacement cost \\
$\mu^{\mathrm{new}}$ & $F \times L$ & Post-replacement mean (default: 0) \\
%% --- CHANGE RF-012: v_new also for rainflow; new b_0 / b_new rows ---
$v^{\mathrm{new}}$ & $F \times L$ & Post-replacement variance
  (Gaussian, Wiener, and Rainflow; default: 0) \\
$b^{0}$ & $F \times L$ & Support upper bound of the initial damage state,
  $D^{0}_{i\ell} \in [0, b^{0}_{i\ell}]$ a.s.\ (Rainflow with
  \texttt{tail\_bound="bernstein"} only; YAML key \texttt{b\_0}).
  Required iff at least one rainflow component uses the Bernstein bound \\
$b^{\mathrm{new}}$ & $F \times L$ & Support upper bound of the post-replacement
  damage state, $D^{\mathrm{new}}_{i\ell} \in [0, b^{\mathrm{new}}_{i\ell}]$ a.s.\
  (Rainflow with \texttt{tail\_bound="bernstein"} only; YAML key
  \texttt{b\_new}).
  Read iff at least one rainflow component uses the Bernstein bound; if absent,
  the default $0$ (deterministic like-new state, matching the
  $\mu^{\mathrm{new}} = v^{\mathrm{new}} = 0$ defaults) is used \\
\bottomrule
\end{longtable}

\subsection{Consistency Checks}
\label{sec:consistency_checks}

The solver must validate all of the following before building the model.
Any violation raises an informative \texttt{ValueError}.

\begin{itemize}[leftmargin=1.5em]
  \item All numerical parameters must be strictly positive where required.
  \item $F > M$.
  \item $\varepsilon \in (0, 0.01]$.
  \item $\mu_0 < \tau$ element-wise (all models).
  \item $\mu^{\mathrm{inc}} < \tau$ element-wise.
  \item $\rho \in (0, 1]$ element-wise.
  %% --- CHANGE MATH-018: add remark about 3*sqrt(v) constraint ---
  \item \textbf{Gaussian:} $\mu_0 \geq 3\sqrt{v_0}$ element-wise;
        $\mu^{\mathrm{inc}} \geq 3\sqrt{v^{\mathrm{inc}}}$ element-wise.

        \emph{Remark.} The $\mu \geq 3\sqrt{v}$ condition ensures that the Gaussian
        degradation distribution has negligible probability of negative values
        (less than $0.13\%$ when $\mu \geq 3\sigma$), maintaining physical validity
        of the non-negative degradation assumption.
  \item \textbf{Gamma:} $\alpha^{\mathrm{inc}} > 0$ element-wise; $\beta > 0$ element-wise.
  \item \textbf{Gaussian:} \texttt{maintenance\_type} must be \texttt{"ARD1"} for
        all Gaussian components (ARA1 is not supported for Gaussian; see
        Section~\ref{sec:dyn_gaussian}).
  \item \textbf{IG:} $\eta > 0$ element-wise; \texttt{maintenance\_type} must be
        \texttt{"ARD1"} for all IG components.
  \item \textbf{Wiener:} $\sigma > 0$ element-wise; $v_0 > 0$ for Wiener components.
  %% --- CHANGE RF-013: rainflow consistency checks ---
  \item \textbf{Rainflow:} \texttt{maintenance\_type} must be \texttt{"ARD1"} for
        all rainflow components (ARA1 is not supported for rainflow, for the same
        reason as for Gaussian: the ARA1 variance semantics are unresolved; see
        Section~\ref{sec:dyn_rainflow} and the Open Questions).
  \item \texttt{tail\_bound} entries must be \texttt{"cantelli"} or
        \texttt{"bernstein"} if specified (single string, or $F \times L$ list);
        entries at non-rainflow positions are ignored with a \texttt{UserWarning}.
        \texttt{"cantelli"} requires no support inputs
        ($b^{\mathrm{inc}}$, $b^{0}$, $b^{\mathrm{new}}$ are not read).
  \item \textbf{Rainflow with \texttt{tail\_bound="bernstein"}:} for every
        applicable component,
        \begin{itemize}[nosep]
          \item $b^{\mathrm{inc}} \geq \mu^{\mathrm{inc}}$ element-wise
                (the mean of a law on $[0, b]$ cannot exceed $b$);
          \item $v^{\mathrm{inc}} \leq \mu^{\mathrm{inc}}
                (b^{\mathrm{inc}} - \mu^{\mathrm{inc}})$ element-wise
                (Bhatia--Davis: this is necessary for \emph{any} law supported on
                $[0, b^{\mathrm{inc}}]$ with the given mean; a violation means the
                measured moments and support are mutually inconsistent);
          \item $b^{0} \geq \mu_0$ element-wise (if $v_0 = 0$,
                $b^{0} = \mu_0$ is admissible); $b^{\mathrm{new}} \geq
                \mu^{\mathrm{new}}$ element-wise; moreover
                $v_0 \leq \mu_0\,(b^{0} - \mu_0)$ and
                $v^{\mathrm{new}} \leq \mu^{\mathrm{new}}\,(b^{\mathrm{new}} -
                \mu^{\mathrm{new}})$ element-wise (Bhatia--Davis, as for the
                increments: a violation means the declared reset moments
                contradict the support claim of assumption A2);
          \item $\tau > 2\kappa\, b_{i\ell}/3$ with $\kappa = \ln(1/\varepsilon)$
                and $b_{i\ell}$ the support constant of
                Section~\ref{sec:dyn_rainflow}: otherwise the Bernstein
                certificate admits no positive variance --- for
                $\tau < 2\kappa b_{i\ell}/3$ it is infeasible even at zero damage
                ($\mu = v = 0$ already requires the margin $2\kappa b_{i\ell}/3$),
                and at equality only the degenerate point $\mu = v = 0$ survives,
                with $V_{\max}^{\mathrm{RF}} = 0$ breaking the big-$M$ rows of
                Section~\ref{sec:dyn_rainflow}.
        \end{itemize}
  %% --- CHANGE MATH-002 option b: require C_D > 0 ---
  \item $C_D > 0$ (required for tightness of the one-sided big-M relaxation;
        see remark in Section~\ref{sec:dyn_gaussian}).
  %% --- CHANGE RF-014: rainflow in allowed model set ---
  \item \texttt{model} entries must all be one of \texttt{"gaussian"},
        \texttt{"inverse\_gaussian"}, \texttt{"wiener"}, \texttt{"gamma"},
        or \texttt{"rainflow"}.
  \item \texttt{maintenance\_type} entries must all be in \texttt{\{"ARA1","ARD1"\}}.
  \item \texttt{penalty\_type} must be \texttt{"inf\_norm"} or \texttt{"quadratic"}
        if specified.
  %% --- CHANGE RF-015: formulation token rename with deprecated alias ---
  \item \texttt{formulation} must be \texttt{"exact"} or \texttt{"lp"} if
        specified. The legacy value \texttt{"socp"} is accepted as a deprecated
        alias: the parser maps it to \texttt{"exact"} and emits a
        \texttt{UserWarning}.
  \item If \texttt{formulation="lp"} and \texttt{penalty\_type="quadratic"}:
        raise \texttt{ValueError} (quadratic penalty requires rotated SOC constraints).
  %% --- CHANGE CS-009: shape validation ---
  \item \textbf{Shape validation.} Before broadcasting, the solver shall validate:
        $\mu^{\mathrm{inc}}$ must be $(F, M, L)$ or $(F, M, L, H)$;
        $v^{\mathrm{inc}}$ must be $(F, M, L)$ or $(F, M, L, H)$;
        $\alpha^{\mathrm{inc}}$ must be $(F, M, L)$ or $(F, M, L, H)$;
        $b^{\mathrm{inc}}$ must be $(F, M, L)$ or $(F, M, L, H)$;
        $C_M, C_R, C_D$ must be scalar, $(F,)$, or $(F, L)$;
        $\alpha_{\mathrm{target}}, \rho, \tau$ must be scalar or $(F, L)$;
        $b^{0}, b^{\mathrm{new}}$ must be scalar or $(F, L)$.
        Raise \texttt{ValueError} with a diagnostic message on shape mismatch.
        Do not rely on NumPy silent broadcasting.
  \item If $H$ is a list: must have exactly two elements; $H_{\min} \geq 1$;
        $H_{\max} > H_{\min}$.
  \item \texttt{n\_workers} must be a positive integer.
  \item If \texttt{warm\_start=True} and \texttt{n\_workers > 1}: raise
        \texttt{ValueError} (warm starting requires sequential execution).
  %% --- CHANGE DOM-026: V_max^user consistency check ---
  %% --- CHANGE RF-016: V_max^user applicability extended to rainflow ---
  \item If $V_{\max}^{\mathrm{user}}$ is specified: must be positive element-wise;
        must satisfy $v_0 \leq V_{\max}^{\mathrm{user}}$ for all applicable components;
        only applicable to Gaussian, Wiener, and Rainflow components.
\end{itemize}

% ================================================================
\section{Decision Variables}
% ================================================================

\begin{longtable}{@{}llp{8cm}@{}}
\toprule
Variable & Shape & Description \\
\midrule
$x_{ijk}$ & $F \times (M{+}1) \times 2H$, binary
  & 1 if train $i$ is assigned to mission $j$ on day $k$
    ($j=0$: maintenance day) \\
$x^m_{i\ell k}$ & $F \times L \times 2H$, binary
  & 1 if component $\ell$ of train $i$ undergoes ARA/ARD maintenance on day $k$ \\
$x^r_{i\ell k}$ & $F \times L \times 2H$, binary
  & 1 if component $\ell$ of train $i$ is replaced on day $k$ \\
$\mu_{i\ell k}$ & $F \times L \times 2H$, non-negative
  & Accumulated mean degradation (all models) \\
%% --- CHANGE RF-017: v tracked by rainflow as well ---
$v_{i\ell k}$ & $F \times L \times 2H$, non-negative
  & Accumulated variance (Gaussian, Wiener, and Rainflow only) \\
$\alpha_{i\ell k}$ & $F \times L \times 2H$, non-negative
  & Accumulated shape parameter (Gamma only);
    $\mu_{i\ell k} = \alpha_{i\ell k}\cdot\beta_{i\ell}$ recovered at post-processing \\
$u$ & scalar, non-negative
  & Damage regularisation variable:
    $u = \max_{i,k}\sum_\ell \mu_{i\ell k}$ (\texttt{inf\_norm}) or
    $u = \max_{i,k}\sum_\ell \mu_{i\ell k}^2$ (\texttt{quadratic}) \\
$z_{i\ell k}$ & $F \times L \times 2H$, non-negative
  & Degradation removed by repair in mean units:
    $z_{i\ell k} = \rho_{i\ell}\,\mu_{i\ell,k-1}\,x^m_{i\ell k}$
    (bilinear; linearised via McCormick envelope) \\
$\mu^{\mathrm{last}}_{i\ell k}$ & $F \times L \times 2H$, non-negative
  & Mean at last maintenance epoch (ARA1 components: Wiener, Gamma) \\
$\alpha^{\mathrm{last}}_{i\ell k}$ & $F \times L \times 2H$, non-negative
  & Shape at last maintenance epoch (ARA1, Gamma only) \\
\bottomrule
\end{longtable}

% ================================================================
\section{Objective Function}
% ================================================================

\begin{align}
\min \quad
  & C_M \sum_{k=1}^{2H} \sum_{i=1}^{F} x_{i0k}
  + \sum_{k=1}^{2H} \sum_{i=1}^{F} \sum_{\ell=1}^{L} C_{R,i\ell}\, z_{i\ell k}
  + \sum_{k=1}^{2H} \sum_{i=1}^{F} \sum_{\ell=1}^{L} C_{\mathrm{rep},i\ell}\, x^r_{i\ell k} \notag \\
  & + C_D\, u
  + C_P \cdot \mathcal{L}_{\mathrm{loop}}
\end{align}

where $z_{i\ell k} = \rho_{i\ell}\,\mu_{i\ell,k-1}\,x^m_{i\ell k}$ is the degradation
removed by repair in mean units (see Section~\ref{sec:repair_proxy}),
$\mathcal{L}_{\mathrm{loop}}$ is the loop-constraint penalty term defined per model
in Section~\ref{sec:loop}, and $u$ is the damage regularisation variable defined by
\texttt{penalty\_type} (see Section~\ref{sec:objective}).

% ================================================================
\section{Constraints}
% ================================================================

\subsection{Assignment Constraints}

Each train performs at most one activity per day:
\begin{equation}
  \sum_{j=0}^{M} x_{ijk} \leq 1 \qquad \forall\, i,k
\end{equation}

Each mission is covered exactly once per day:
\begin{equation}
  \sum_{i=1}^{F} x_{ijk} = 1 \qquad \forall\, j \geq 1,\; k
\end{equation}

Together these constraints guarantee that exactly $M$ trains are on missions at every
step $k$, and each train performs at most one activity per day. The remaining $F - M$
trains may be on maintenance or idle. No additional fleet-availability constraint is needed.

\subsection{Maintenance / Replacement Compatibility}

Maintenance and replacement may only occur on a maintenance day ($x_{i0k} = 1$):
\begin{equation}
  x^m_{i\ell k} \leq x_{i0k} \qquad \forall\, i,\ell,k
\end{equation}
\begin{equation}
  x^r_{i\ell k} \leq x_{i0k} \qquad \forall\, i,\ell,k
\end{equation}

A component cannot be both repaired and replaced on the same day:
\begin{equation}
  x^m_{i\ell k} + x^r_{i\ell k} \leq 1 \qquad \forall\, i,\ell,k
\end{equation}

\subsection{State Dynamics --- Gaussian Model}
\label{sec:dyn_gaussian}

For components $(i,\ell)$ with \texttt{model} $=$ \texttt{gaussian}.
\textbf{Only \texttt{ARD1} is supported} (consistency check must enforce this).
Both $\mu_{i\ell k}$ and $v_{i\ell k}$ are tracked.

\textbf{Notation.}
$\delta\mu_{i\ell k} = \sum_{j=1}^{M} x_{ijk}\,\mu^{\mathrm{inc}}_{ij\ell k}$,
$\delta v_{i\ell k} = \sum_{j=1}^{M} x_{ijk}\,v^{\mathrm{inc}}_{ij\ell k}$
(mission-induced increments; periodic extension for $k>H$).
Big-M values: $M_\mu = \tau$, $M_v = V_{\max} = \tau^2/(\Phi^{-1}(1-\varepsilon))^2$.

%% --- CHANGE MATH-003: big-M validity remark ---
\emph{Remark.} The big-M value $M_\mu = \tau$ is valid because the reliability
constraint ensures $\mu_{i\ell k} \leq \tau$ for all $i, \ell, k$.
This bound is also consistent with the physical interpretation that degradation
beyond $\tau$ triggers failure.

\subsubsection{ARD1 Gaussian}

\textbf{Mean dynamics --- exact form:}
\begin{equation}
  \mu_{i\ell k} = \begin{cases}
    \mu_{i\ell,k-1} + \delta\mu_{i\ell k} & \text{if } x^m_{i\ell k} = 0 \\
    (1-\rho_{i\ell})\,\mu_{i\ell,k-1} & \text{if } x^m_{i\ell k} = 1
  \end{cases}
\end{equation}

%% --- CHANGE MATH-019: add -tau*x^r term to accumulation constraints ---
%% --- CHANGE DOM-012: tightness remark ---
\textbf{Mean dynamics --- linearised (big-M):}
\begin{align}
  \mu_{i\ell k} &\geq \mu_{i\ell,k-1} + \delta\mu_{i\ell k} - \tau\,x^m_{i\ell k}
    - \tau\,x^r_{i\ell k}
    \label{eq:g_ard_mu1} \\
  \mu_{i\ell k} &\geq (1-\rho_{i\ell})\,\mu_{i\ell,k-1} - \tau(1-x^m_{i\ell k})
    \label{eq:g_ard_mu2}
\end{align}

%% --- CHANGE MATH-002: cross-ref reliability for Eq. 9 vacuousness ---
\emph{Remark (vacuousness when $x^m=0$).}
When $x^m_{i\ell k} = 0$, the accumulation constraint reduces to
$\mu_{i\ell k} \geq \mu_{i\ell,k-1} + \delta\mu_{i\ell k}$, which is the
only active bound. The non-trivial control on~$\mu$ in non-maintenance periods
comes from the reliability constraint (Section~\ref{sec:dyn_gaussian}, Reliability),
which enforces $\mu_{i\ell k} \leq \tau$ (or tighter, depending on $\alpha$).

%% --- CHANGE MATH-004+DEV-001: refined three-mechanism tightness argument ---
\emph{Remark (tightness of big-M relaxation at optimality).}
The one-sided relaxation $\mu^+_{i\ell k} \leq \mu^-_{i\ell k}$ (instead of equality
when $x^m_{i\ell k}=0$) is tight at any optimal solution, through three complementary
mechanisms:
\begin{enumerate}[nosep]
  \item $C_D \cdot u$: at the arg-max $(i,k)$ pair where the damage penalty is binding,
        inflating $\mu$ directly increases cost.
  \item $C_R \cdot z_{i\ell k}$: for every subsequent maintenance step $k'>k$,
        forward propagation of an inflated $\mu_{i\ell k}$ through the degradation
        dynamics increases the repair cost $z_{i\ell k'}$.
  \item $C_P \cdot \mathcal{L}_{\mathrm{loop}}$: at half-horizon boundaries
        ($k=H$ and $k=2H$), any slack in $\mu$ worsens the loop penalty.
\end{enumerate}
%% --- CHANGE RF-018: SOC wording updated (exact form is not SOC) ---
For the variance variable~$v$, tightness is enforced by the reliability constraint:
an inflated $v$ tightens the upper bound on $\mu$ via the exact or LP reliability
form, creating an indirect cost penalty.

%% --- CHANGE MATH-003: ARD1 deterministic variance scaling remark ---
\emph{Remark (ARD1 variance scaling).}
Under ARD1, the post-repair variance follows
$v^+_{i\ell k} = (1-\rho_{i\ell})^2\,v^-_{i\ell k}$. This is a direct algebraic
consequence of the degradation-level reduction $D^+ = (1-\rho)D^-$ applied to a
Gaussian random variable: $\mathrm{Var}[(1-\rho)D] = (1-\rho)^2\,\mathrm{Var}[D]$.
No additional modelling assumption is required.

\textbf{Variance dynamics --- exact form:}
\begin{equation}
  v_{i\ell k} = \begin{cases}
    v_{i\ell,k-1} + \delta v_{i\ell k} & \text{if } x^m_{i\ell k} = 0 \\
    (1-\rho_{i\ell})^2\,v_{i\ell,k-1} & \text{if } x^m_{i\ell k} = 1
  \end{cases}
\end{equation}

%% --- CHANGE MATH-019: add -V_max*x^r term to variance accumulation ---
\textbf{Variance dynamics --- linearised (big-M):}
\begin{align}
  v_{i\ell k} &\geq v_{i\ell,k-1} + \delta v_{i\ell k} - V_{\max}\,x^m_{i\ell k}
    - V_{\max}\,x^r_{i\ell k} \\
  v_{i\ell k} &\geq (1-\rho_{i\ell})^2\,v_{i\ell,k-1} - V_{\max}(1-x^m_{i\ell k})
\end{align}

\textbf{Replacement --- exact form:}
\begin{equation}
  (\mu_{i\ell k},\, v_{i\ell k}) = (\mu^{\mathrm{new}}_{i\ell},\, v^{\mathrm{new}}_{i\ell})
  \quad \text{if } x^r_{i\ell k} = 1
  \label{eq:g_rep}
\end{equation}

\textbf{Replacement --- linearised (big-M):}
\begin{align}
  \mu_{i\ell k} &\leq \mu^{\mathrm{new}}_{i\ell} + \tau(1-x^r_{i\ell k}), \quad
  \mu_{i\ell k} \geq \mu^{\mathrm{new}}_{i\ell} - \tau(1-x^r_{i\ell k}) \\
  v_{i\ell k} &\leq v^{\mathrm{new}}_{i\ell} + V_{\max}(1-x^r_{i\ell k}), \quad
  v_{i\ell k} \geq v^{\mathrm{new}}_{i\ell} - V_{\max}(1-x^r_{i\ell k})
\end{align}

%% --- CHANGE DOM-009: ARA1 Gaussian removed ---
%% ARA1 is not supported for Gaussian components. The ARA1 variance extension
%% (v+ = v- - rho*(v- - v^last)) has no established basis in the maintenance
%% literature (Mercier & Castro 2019 define ARA1 for Gamma via virtual age,
%% not explicit variance tracking). Gaussian uses ARD1 only, which provides
%% deterministic (1-rho)^2 variance scaling as shown above.

\subsubsection{Reliability Constraint (Gaussian)}

%% --- CHANGE MATH-006: mu <= tau remark ---
%% --- CHANGE RF-019: formulation token socp->exact; rotated-SOC mislabel corrected ---
\textbf{Exact form (nonconvex quadratic, \texttt{formulation="exact"}):}
\begin{equation}
  \mu_{i\ell k} + \Phi^{-1}(1-\varepsilon)\,\sqrt{v_{i\ell k}} \leq \tau
  \qquad \forall\, i,\ell,k
\end{equation}
Equivalently, squaring both sides (valid since $\tau - \mu_{i\ell k} \geq 0$ by feasibility
and $\Phi^{-1}(1-\varepsilon) > 0$ since $\varepsilon \leq 0.01$):
\begin{equation}
  (\Phi^{-1}(1-\varepsilon))^2\,v_{i\ell k} \leq (\tau - \mu_{i\ell k})^2
\end{equation}
This squared form is exact but \textbf{not} a rotated second-order cone: since
the variance $v_{i\ell k}$ is itself an (affine) decision quantity, the
constraint reads $(\text{affine}) \leq (\text{affine})^2$, and the
$\mu_{i\ell k}^2$ term enters with the wrong sign for convexity---the feasible
set is the complement of a convex set. Concretely,
$(\mu_{i\ell k}, v_{i\ell k}) = (\tau,\, 0)$ and
$(0,\, \tau^2/(\Phi^{-1}(1-\varepsilon))^2)$ are both feasible while their
midpoint $(\tau/2,\, \tau^2/(2(\Phi^{-1}(1-\varepsilon))^2))$ violates the
constraint ($\tau^2/2 > \tau^2/4$). Gurobi handles the constraint through its
nonconvex quadratic machinery (parameter \texttt{NonConvex=2} where not the
default); the constraint remains exact---no approximation is introduced. The
legacy value \texttt{formulation="socp"} is accepted as a deprecated alias for
\texttt{"exact"} (Section~\ref{sec:input_scalars}).

\emph{Remark.} The bound $\mu_{i\ell k} \leq \tau$ is guaranteed by $\mu_0 < \tau$,
the non-negativity of degradation increments, and the dynamics constraints;
it does not rely on the reliability constraint itself.

\textbf{LP approximation (\texttt{formulation="lp"}):}
Expanding the right-hand side and dropping the non-negative term $\mu_{i\ell k}^2$:
\begin{equation}
  (\Phi^{-1}(1-\varepsilon))^2\,v_{i\ell k} + 2\tau\,\mu_{i\ell k} \leq \tau^2
  \qquad \forall\, i,\ell,k
  \label{eq:g_lp_reliability}
\end{equation}
%% --- CHANGE RF-020: SOC wording updated in LP derivation and conservatism remark ---
This is an \textbf{inner approximation}: any $(v_{i\ell k}, \mu_{i\ell k})$ satisfying
Eq.~\eqref{eq:g_lp_reliability} also satisfies the exact squared constraint,
since $(\tau - \mu_{i\ell k})^2 = \tau^2 - 2\tau\,\mu_{i\ell k} + \mu_{i\ell k}^2
\geq \tau^2 - 2\tau\,\mu_{i\ell k}$.
The approximation error is $\mu_{i\ell k}^2$, which is small when $\mu_{i\ell k} \ll \tau$
(the typical operating regime) and vanishes when $\tau = 1$ and $\mu_{i\ell k} \to 0$.
No additional parameters are required.

%% --- CHANGE DEV-007: quantify LP approximation conservatism ---
\emph{Remark (LP approximation conservatism).}
The LP inner approximation replaces the exact quadratic constraint with a tangent
hyperplane, yielding a feasible region that is a strict subset of the exact
feasible region.
This means the LP relaxation may exclude some exact-feasible points, potentially
increasing the objective value. The approximation is tightest when $v_{i\ell k}$
is near the linearisation point and loosest at the boundaries of the variance domain.
When solution quality is critical, the exact formulation should be preferred;
the LP approximation is intended for solvers without nonconvex-quadratic support.

%% --- CHANGE DOM-026: user-defined variance upper bound constraint ---
\textbf{User-defined variance upper bound (optional):}
If $V_{\max}^{\mathrm{user}}$ is specified for component $(i,\ell)$:
\begin{equation}
  v_{i\ell k} \leq V_{\max,i\ell}^{\mathrm{user}} \qquad \forall\, k
\end{equation}

% ----------------------------------------------------------------
\subsection{State Dynamics --- Inverse Gaussian Model}
\label{sec:dyn_ig}
% ----------------------------------------------------------------

For components $(i,\ell)$ with \texttt{model} $=$ \texttt{inverse\_gaussian}.
\textbf{Only \texttt{ARD1} is supported} (consistency check must enforce this).
Only $\mu_{i\ell k}$ is tracked; $\lambda_{i\ell k} = \mu_{i\ell k}/\eta_{i\ell}$
is recovered analytically at post-processing. Big-M: $M_\mu = \tau$.

\textbf{Notation.} $\delta\mu_{i\ell k} = \sum_{j=1}^{M} x_{ijk}\,\mu^{\mathrm{inc}}_{ij\ell k}$
(periodic extension for $k > H$).

\subsubsection{ARD1 Inverse Gaussian}

\textbf{Mean dynamics --- exact form:}
\begin{equation}
  \mu_{i\ell k} = \begin{cases}
    \mu_{i\ell,k-1} + \delta\mu_{i\ell k} & \text{if } x^m_{i\ell k} = 0 \\
    (1-\rho_{i\ell})\,\mu_{i\ell,k-1} & \text{if } x^m_{i\ell k} = 1
  \end{cases}
\end{equation}

%% --- CHANGE MATH-019: add -tau*x^r term ---
\textbf{Mean dynamics --- linearised (big-M):}
\begin{align}
  \mu_{i\ell k} &\geq \mu_{i\ell,k-1} + \delta\mu_{i\ell k} - \tau\,x^m_{i\ell k}
    - \tau\,x^r_{i\ell k}
    \label{eq:ig_mu1} \\
  \mu_{i\ell k} &\geq (1-\rho_{i\ell})\,\mu_{i\ell,k-1} - \tau(1-x^m_{i\ell k})
    \label{eq:ig_mu2}
\end{align}

\textbf{Replacement --- exact form:}
\begin{equation}
  \mu_{i\ell k} = \mu^{\mathrm{new}}_{i\ell} \quad \text{if } x^r_{i\ell k} = 1
\end{equation}

\textbf{Replacement --- linearised (big-M):}
\begin{align}
  \mu_{i\ell k} &\leq \mu^{\mathrm{new}}_{i\ell} + \tau(1-x^r_{i\ell k}) \\
  \mu_{i\ell k} &\geq \mu^{\mathrm{new}}_{i\ell} - \tau(1-x^r_{i\ell k})
\end{align}

\subsubsection{Reliability Constraint (IG)}

\textbf{Exact form (nonlinear, offline-precomputed):}
The constraint $\Pr(D_{i\ell k} > \tau) \leq \varepsilon$ with
$D \sim \mathrm{IG}(\mu_{i\ell k},\,\mu_{i\ell k}/\eta_{i\ell})$
is monotone increasing in $\mu_{i\ell k}$, so it reduces to:
\begin{equation}
  \mu_{i\ell k} \leq \bar{\mu}_{i\ell} \qquad \forall\, i,\ell,k
\end{equation}
where $\bar{\mu}_{i\ell}$ is the unique root of:
\begin{equation}
  \Phi\!\left(\sqrt{\frac{\bar\mu}{\eta_{i\ell}\,\tau}}
    \left(\frac{\tau}{\bar\mu} - 1\right)\right)
  + e^{2/\eta_{i\ell}}\,
  \Phi\!\left(-\sqrt{\frac{\bar\mu}{\eta_{i\ell}\,\tau}}
    \left(\frac{\tau}{\bar\mu} + 1\right)\right)
  = 1 - \varepsilon
  \label{eq:ig_mubar}
\end{equation}

%% --- CHANGE MATH-008: monotonicity argument for IG reliability ---
\emph{Remark (monotonicity).}
Monotonicity of $\Pr(D > \tau)$ in $\mu$ follows from the likelihood-ratio ordering
of IG distributions at fixed $\eta$: if $\mu_1 \leq \mu_2$, then
$\mathrm{IG}(\mu_1, \mu_1/\eta) \leq_{\mathrm{lr}} \mathrm{IG}(\mu_2, \mu_2/\eta)$,
which implies $\Pr(D_1 > \tau) \leq \Pr(D_2 > \tau)$.

%% --- CHANGE CS-010: provide log-sum-exp form and scipy function ---
\textbf{Linearised form:} since $\bar{\mu}_{i\ell}$ depends only on inputs
$(\eta_{i\ell}, \tau, \varepsilon)$, it is computed \textbf{offline} before the Gurobi
model is built. The constraint entering Gurobi is the linear bound
$\mu_{i\ell k} \leq \bar{\mu}_{i\ell}$.

\emph{Numerical implementation.}
The IG CDF involves a term $\exp(2\lambda/\mu) = \exp(2/\eta)$ that can overflow.
Compute in log-space: $\log p_2 = 2/\eta_{i\ell} + \log\Phi(\cdot)$, then use
\texttt{logsumexp} to combine. In \texttt{scipy}:
use \texttt{scipy.stats.invgauss.cdf(x, mu=mu/lam, scale=lam)} with
\texttt{scipy.optimize.brentq} for root-finding on
Eq.~\eqref{eq:ig_mubar}.

\subsubsection{Loop Constraint (IG, hard + soft)}

%% --- CHANGE MATH-001+DEV-005+DOM-012: replace TODO with verified proposition ---
\emph{Proposition (IG FSD via MLR at constant $\eta$).}
Let $D_1 \sim \mathrm{IG}(\mu_1, \mu_1/\eta)$ and
$D_2 \sim \mathrm{IG}(\mu_2, \mu_2/\eta)$ with $\mu_1 \leq \mu_2$ and $\eta > 0$
constant. Then $D_1 \leq_{\mathrm{FSD}} D_2$.

\emph{Proof.}
The density of $\mathrm{IG}(\mu, \lambda)$ with $\lambda = \mu/\eta$ is
$f(x; \mu) = \sqrt{\mu/(2\pi\eta x^3)}\,
\exp\!\bigl(-{(x-\mu)^2}/{(2\eta x)}\bigr)$ for $x > 0$.
The log-likelihood ratio satisfies:
\[
  \frac{\partial}{\partial x} \log \frac{f(x;\mu_2)}{f(x;\mu_1)}
  = \frac{\mu_2 - \mu_1}{2\eta}\!\left(\frac{1}{\mu_1 \mu_2} + \frac{1}{x^2}\right)
  \geq 0
  \quad \forall\, x > 0,
\]
establishing the monotone likelihood ratio (MLR) property. By the standard chain
MLR $\Rightarrow$ HR $\Rightarrow$ FSD, the result follows.\hfill$\square$

Since ARD1 repair scales both $\mu$ and $\lambda = \mu/\eta$ by the same factor
$(1-\rho)$, the ratio $\eta$ is preserved identically. Therefore FSD reduces to the
mean-only condition:
FSD reduces to a mean-only condition:
\begin{equation}
  \mu_{i\ell,2H} \leq \mu_{i\ell,H}
  \qquad \forall\,(i,\ell)\text{ s.t. model} = \mathrm{IG}
\end{equation}
\begin{equation}
  \mathcal{L}^{\mathrm{IG}}_{\mathrm{loop}} =
    \sum_{(i,\ell)\,:\,\mathrm{IG}} \bigl(\mu_{i\ell H} - \mu_{i\ell 2H}\bigr)
\end{equation}

% ----------------------------------------------------------------
\subsection{State Dynamics --- Wiener Model}
\label{sec:dyn_wiener}
% ----------------------------------------------------------------

For components $(i,\ell)$ with \texttt{model} $=$ \texttt{wiener}.
Supports both \texttt{ARD1} and \texttt{ARA1} maintenance types.
Both $\mu_{i\ell k}$ and $v_{i\ell k}$ are tracked.
\textbf{Imperfect repair acts on the mean only}; variance is unaffected by repair
and can only be reset by replacement.
Variance accumulates during missions only (not on maintenance or idle days).

\textbf{Notation.}
$\delta\mu_{i\ell k} = \sum_{j=1}^{M} x_{ijk}\,\mu^{\mathrm{inc}}_{ij\ell k}$,
$\delta v_{i\ell k} = \sigma^2_{i\ell}\sum_{j=1}^{M} x_{ijk}$
(mission-only variance increment; $\sigma_{i\ell}$ is the diffusion constant).
Big-M values: $M_\mu = \tau$, $M_v = V_{\max} = \tau^2/(\Phi^{-1}(1-\varepsilon))^2$.

%% --- CHANGE DOM-006 option a: justify mission-independent Wiener variance ---
\emph{Remark.} The Wiener variance increment $\delta v_{i\ell k} = \sigma^2_{i\ell}\sum_j x_{ijk}$
is mission-independent (the same $\sigma_{i\ell}$ applies regardless of which mission
$j$ is performed). This reflects the physical interpretation that $\sigma$ characterises
material microstructure and diffusion properties, which are intrinsic to the component
and independent of the route (mission) performed.

\subsubsection{ARD1 Wiener}

\textbf{Mean dynamics --- exact form:}
\begin{equation}
  \mu_{i\ell k} = \begin{cases}
    \mu_{i\ell,k-1} + \delta\mu_{i\ell k} & \text{if } x^m_{i\ell k} = 0 \\
    (1-\rho_{i\ell})\,\mu_{i\ell,k-1} & \text{if } x^m_{i\ell k} = 1
  \end{cases}
\end{equation}

%% --- CHANGE MATH-019: add -tau*x^r term ---
\textbf{Mean dynamics --- linearised (big-M):}
\begin{align}
  \mu_{i\ell k} &\geq \mu_{i\ell,k-1} + \delta\mu_{i\ell k} - \tau\,x^m_{i\ell k}
    - \tau\,x^r_{i\ell k} \\
  \mu_{i\ell k} &\geq (1-\rho_{i\ell})\,\mu_{i\ell,k-1} - \tau(1-x^m_{i\ell k})
\end{align}

\textbf{Variance dynamics --- exact form} (repair-independent):
\begin{equation}
  v_{i\ell k} = v_{i\ell,k-1} + \sigma^2_{i\ell}\sum_{j=1}^{M} x_{ijk}
  \quad \text{(if } x^r_{i\ell k} = 0\text{)}
\end{equation}

%% --- CHANGE MATH-011: tightness remark for Wiener variance ---
\textbf{Variance dynamics --- linearised (big-M):}
\begin{equation}
  v_{i\ell k} \geq v_{i\ell,k-1} + \delta v_{i\ell k} - V_{\max}\,x^r_{i\ell k}
\end{equation}
The $x^r_{i\ell k}$ term allows replacement to reset variance; $x^m_{i\ell k}$
does not appear since repair does not affect variance.

%% --- CHANGE MATH-009: Wiener variance tightness via reliability ---
\emph{Remark (Wiener variance tightness).}
Unlike the mean, whose tightness follows from the three cost mechanisms
(Section~\ref{sec:dyn_gaussian}, tightness remark), the Wiener variance
$v_{i\ell k} = v_0 + k\sigma^2_{i\ell}$ is not directly penalised in the objective.
Tightness of the variance relaxation is enforced indirectly by the reliability
constraint: an inflated $v_{i\ell k}$ reduces the effective failure threshold via
$\mu_{i\ell k} \leq \tau - \Phi^{-1}(1-\varepsilon)\sqrt{v_{i\ell k}}$, which
tightens the feasible region for $\mu$ and increases cost. This mechanism is active
whenever the reliability constraint is binding.

\textbf{Replacement --- exact form:}
\begin{equation}
  (\mu_{i\ell k},\, v_{i\ell k}) = (\mu^{\mathrm{new}}_{i\ell},\, v^{\mathrm{new}}_{i\ell})
  \quad \text{if } x^r_{i\ell k} = 1
\end{equation}

\textbf{Replacement --- linearised (big-M):}
\begin{align}
  \mu_{i\ell k} &\leq \mu^{\mathrm{new}}_{i\ell} + \tau(1-x^r_{i\ell k}), \quad
  \mu_{i\ell k} \geq \mu^{\mathrm{new}}_{i\ell} - \tau(1-x^r_{i\ell k}) \\
  v_{i\ell k} &\leq v^{\mathrm{new}}_{i\ell} + V_{\max}(1-x^r_{i\ell k}), \quad
  v_{i\ell k} \geq v^{\mathrm{new}}_{i\ell} - V_{\max}(1-x^r_{i\ell k})
\end{align}

\subsubsection{ARA1 Wiener}

Requires auxiliary variable $\mu^{\mathrm{last}}_{i\ell k}$ only
(no $v^{\mathrm{last}}_{i\ell k}$ since repair does not affect variance).
Initial value: $\mu^{\mathrm{last}}_{i\ell,0} = \mu_{0,i\ell}$.

\textbf{Auxiliary update --- exact (bilinear) form:}
\begin{equation}
  \mu^{\mathrm{last}}_{i\ell k} = \begin{cases}
    \mu_{i\ell,k-1} & \text{if } x^m_{i\ell k} = 1 \\
    \mu^{\mathrm{last}}_{i\ell,k-1} & \text{if } x^m_{i\ell k} = 0
  \end{cases}
\end{equation}

\textbf{Auxiliary update --- linearised (big-M, $M_\mu = \tau$):}
\begin{align}
  \mu^{\mathrm{last}}_{i\ell k} &\geq \mu_{i\ell,k-1} - \tau(1-x^m_{i\ell k}), \quad
  \mu^{\mathrm{last}}_{i\ell k} \leq \mu_{i\ell,k-1} + \tau(1-x^m_{i\ell k}) \\
  \mu^{\mathrm{last}}_{i\ell k} &\geq \mu^{\mathrm{last}}_{i\ell,k-1} - \tau\,x^m_{i\ell k}, \quad
  \mu^{\mathrm{last}}_{i\ell k} \leq \mu^{\mathrm{last}}_{i\ell,k-1} + \tau\,x^m_{i\ell k}
\end{align}

\textbf{Mean dynamics --- exact form:}
\begin{equation}
  \mu_{i\ell k} = \begin{cases}
    \mu_{i\ell,k-1} + \delta\mu_{i\ell k} & \text{if } x^m_{i\ell k} = 0 \\
    \mu_{i\ell,k-1} - \rho_{i\ell}(\mu_{i\ell,k-1} - \mu^{\mathrm{last}}_{i\ell,k-1})
      & \text{if } x^m_{i\ell k} = 1
  \end{cases}
\end{equation}

%% --- CHANGE MATH-019: add -tau*x^r term ---
%% --- CHANGE MATH-021: cross-reference to tightness argument ---
\textbf{Mean dynamics --- linearised (big-M):}
\begin{align}
  \mu_{i\ell k} &\geq \mu_{i\ell,k-1} + \delta\mu_{i\ell k} - \tau\,x^m_{i\ell k}
    - \tau\,x^r_{i\ell k} \\
  \mu_{i\ell k} &\geq \mu_{i\ell,k-1}
    - \rho_{i\ell}(\mu_{i\ell,k-1} - \mu^{\mathrm{last}}_{i\ell,k-1})
    - \tau(1-x^m_{i\ell k})
\end{align}

\emph{Remark.} As with all one-sided dynamics linearisations in this specification,
tightness at optimality is guaranteed by the $C_D > 0$ requirement
(see the tightness remark in Section~\ref{sec:dyn_gaussian}).

\textbf{Variance dynamics and replacement:} identical to ARD1 Wiener.

\subsubsection{Reliability Constraint (Wiener)}

%% --- CHANGE RF-021: formulation token socp->exact; rotated-SOC mislabel corrected (Wiener) ---
\textbf{Exact form (nonconvex quadratic, \texttt{formulation="exact"}):} identical to Gaussian,
\begin{equation}
  \mu_{i\ell k} + \Phi^{-1}(1-\varepsilon)\,\sqrt{v_{i\ell k}} \leq \tau
  \qquad \forall\, i,\ell,k
\end{equation}
Equivalently: $(\Phi^{-1}(1-\varepsilon))^2\,v_{i\ell k} \leq (\tau - \mu_{i\ell k})^2$,
an exact \emph{nonconvex} quadratic constraint (the $\mu_{i\ell k}^2$ term enters
with the wrong sign for convexity; see the corrected discussion in
Section~\ref{sec:dyn_gaussian}). Handled by Gurobi's nonconvex quadratic support
(\texttt{NonConvex=2} where not the default); no approximation is introduced.

\textbf{LP approximation (\texttt{formulation="lp"}):} identical form to Gaussian,
\begin{equation}
  (\Phi^{-1}(1-\varepsilon))^2\,v_{i\ell k} + 2\tau\,\mu_{i\ell k} \leq \tau^2
  \qquad \forall\, i,\ell,k
  \label{eq:w_lp_reliability}
\end{equation}
Same inner approximation argument applies; approximation error is $\mu_{i\ell k}^2$.

%% --- CHANGE DOM-026: user-defined variance upper bound constraint for Wiener ---
\textbf{User-defined variance upper bound (optional):}
If $V_{\max}^{\mathrm{user}}$ is specified for Wiener component $(i,\ell)$:
\begin{equation}
  v_{i\ell k} \leq V_{\max,i\ell}^{\mathrm{user}} \qquad \forall\, k
\end{equation}

\subsubsection{Loop Constraint (Wiener, hard + soft)}

%% --- CHANGE RF-022: A-02 accepted review item applied (FSD wording corrected) ---
The loop constraint bounds both tracked moments,
$\mu_{i\ell,2H} \leq \mu_{i\ell,H}$ and $v_{i\ell,2H} \leq v_{i\ell,H}$, to
guarantee perpetual safety (see Section~\ref{sec:loop}). For Gaussian marginals
with $\mu < \tau$, bounding both moments implies the failure-probability
comparison $\Pr(D_{2H} > \tau) \leq \Pr(D_H > \tau)$ at any threshold $\tau$,
since $\Phi\bigl((\mu - \tau)/\sqrt{v}\bigr)$ is increasing in $\mu$ and, for
$\mu < \tau$, increasing in $v$. It is \emph{not} first-order stochastic
dominance: Gaussian CDFs with different variances cross, so no FSD relation
holds between them. Since variance can only be reset by replacement,
satisfying $v_{i\ell,2H} \leq v_{i\ell,H}$ may require replacement actions
within $[H+1,2H]$; infeasibility is possible if $\sigma_{i\ell}$ is large relative
to $v_{i\ell,H} - v^{\mathrm{new}}_{i\ell}$.

%% --- CHANGE MATH-015 / DOM-019 / DEV-007 (REJECTED): add clarifying remark ---
\emph{Remark.} This constraint intentionally enforces at least one replacement
in $[H{+}1,2H]$ for Wiener components, since imperfect repair cannot reduce variance.

Hard constraints:
\begin{equation}
  \mu_{i\ell,2H} \leq \mu_{i\ell,H}, \quad
  v_{i\ell,2H} \leq v_{i\ell,H}
  \qquad \forall\,(i,\ell)\text{ s.t. model} = \mathrm{Wiener}
\end{equation}
Loop penalty contribution:
\begin{equation}
  \mathcal{L}^{\mathrm{Wiener}}_{\mathrm{loop}} =
    \sum_{(i,\ell)\,:\,\mathrm{Wiener}}
    \bigl[(\mu_{i\ell H} - \mu_{i\ell,2H}) + (v_{i\ell H} - v_{i\ell,2H})\bigr]
\end{equation}

% ----------------------------------------------------------------
\subsection{State Dynamics --- Gamma Model}
\label{sec:dyn_gamma}
% ----------------------------------------------------------------

For components $(i,\ell)$ with \texttt{model} $=$ \texttt{gamma}.
Supports both \texttt{ARD1} and \texttt{ARA1} maintenance types.

\textbf{State variable.} $\alpha_{i\ell k}$ is the accumulated Gamma shape parameter.
Initial value: $\alpha_{i\ell,0} = \mu_{0,i\ell}/\beta_{i\ell}$.
Post-replacement value: $\alpha^{\mathrm{new}}_{i\ell} = \mu^{\mathrm{new}}_{i\ell}/\beta_{i\ell}$.
Both are derived from user-specified inputs and $\beta_{i\ell}$.
The mean is recovered at post-processing as $\mu_{i\ell k} = \alpha_{i\ell k}\cdot\beta_{i\ell}$.

\textbf{Big-M.} The tightest valid big-M for $\alpha_{i\ell k}$ is:
\begin{equation}
  \bar{\alpha}_{i\ell} = \frac{\tau}{\beta_{i\ell}}
\end{equation}
since $\alpha_{i\ell k} \leq \tau/\beta_{i\ell}$ is implied by $\mu_{i\ell k} < \tau$ and $\beta_{i\ell} > 0$.

\textbf{Notation.} Let $\delta\alpha_{i\ell k} = \sum_{j=1}^{M} x_{ijk}\,\alpha^{\mathrm{inc}}_{ij\ell k}$
be the mission-induced shape increment at step $k$ (periodic extension for $k > H$).

\subsubsection{ARD1 Gamma}

\textbf{Shape dynamics --- exact form:}
\begin{equation}
  \alpha_{i\ell k} = \begin{cases}
    \alpha_{i\ell,k-1} + \delta\alpha_{i\ell k} & \text{if } x^m_{i\ell k} = 0 \\
    (1-\rho_{i\ell})\,\alpha_{i\ell,k-1} & \text{if } x^m_{i\ell k} = 1
  \end{cases}
\end{equation}

%% --- CHANGE MATH-019: add -bar_alpha*x^r term ---
\textbf{Shape dynamics --- linearised (big-M):}
\begin{align}
  \alpha_{i\ell k} &\geq \alpha_{i\ell,k-1} + \delta\alpha_{i\ell k}
    - \bar{\alpha}_{i\ell}\,x^m_{i\ell k}
    - \bar{\alpha}_{i\ell}\,x^r_{i\ell k} \\
  \alpha_{i\ell k} &\geq (1-\rho_{i\ell})\,\alpha_{i\ell,k-1}
    - \bar{\alpha}_{i\ell}(1 - x^m_{i\ell k})
\end{align}

\textbf{Replacement --- exact form:}
\begin{equation}
  \alpha_{i\ell k} = \alpha^{\mathrm{new}}_{i\ell} \quad \text{if } x^r_{i\ell k} = 1
\end{equation}

\textbf{Replacement --- linearised (big-M):}
\begin{align}
  \alpha_{i\ell k} &\leq \alpha^{\mathrm{new}}_{i\ell}
    + \bar{\alpha}_{i\ell}(1 - x^r_{i\ell k}) \\
  \alpha_{i\ell k} &\geq \alpha^{\mathrm{new}}_{i\ell}
    - \bar{\alpha}_{i\ell}(1 - x^r_{i\ell k})
\end{align}

\subsubsection{ARA1 Gamma}

Requires auxiliary variable $\alpha^{\mathrm{last}}_{i\ell k}$, tracking the shape
immediately after the most recent maintenance event.
Initial value: $\alpha^{\mathrm{last}}_{i\ell,0} = \alpha_{i\ell,0} = \mu_{0,i\ell}/\beta_{i\ell}$.

\textbf{Auxiliary update --- exact (bilinear) form:}
\begin{equation}
  \alpha^{\mathrm{last}}_{i\ell k} = \begin{cases}
    \alpha_{i\ell,k-1} & \text{if } x^m_{i\ell k} = 1 \\
    \alpha^{\mathrm{last}}_{i\ell,k-1} & \text{if } x^m_{i\ell k} = 0
  \end{cases}
\end{equation}

\textbf{Auxiliary update --- linearised (big-M $= \bar{\alpha}_{i\ell}$):}
\begin{align}
  \alpha^{\mathrm{last}}_{i\ell k} &\geq \alpha_{i\ell,k-1}
    - \bar{\alpha}_{i\ell}(1 - x^m_{i\ell k}), \quad
  \alpha^{\mathrm{last}}_{i\ell k} \leq \alpha_{i\ell,k-1}
    + \bar{\alpha}_{i\ell}(1 - x^m_{i\ell k}) \\
  \alpha^{\mathrm{last}}_{i\ell k} &\geq \alpha^{\mathrm{last}}_{i\ell,k-1}
    - \bar{\alpha}_{i\ell}\,x^m_{i\ell k}, \quad
  \alpha^{\mathrm{last}}_{i\ell k} \leq \alpha^{\mathrm{last}}_{i\ell,k-1}
    + \bar{\alpha}_{i\ell}\,x^m_{i\ell k}
\end{align}

\textbf{Shape dynamics --- exact form:}
\begin{equation}
  \alpha_{i\ell k} = \begin{cases}
    \alpha_{i\ell,k-1} + \delta\alpha_{i\ell k} & \text{if } x^m_{i\ell k} = 0 \\
    \alpha_{i\ell,k-1} - \rho_{i\ell}(\alpha_{i\ell,k-1} - \alpha^{\mathrm{last}}_{i\ell,k-1})
      & \text{if } x^m_{i\ell k} = 1
  \end{cases}
\end{equation}

%% --- CHANGE MATH-019: add -bar_alpha*x^r term ---
\textbf{Shape dynamics --- linearised (big-M):}
\begin{align}
  \alpha_{i\ell k} &\geq \alpha_{i\ell,k-1} + \delta\alpha_{i\ell k}
    - \bar{\alpha}_{i\ell}\,x^m_{i\ell k}
    - \bar{\alpha}_{i\ell}\,x^r_{i\ell k} \\
  \alpha_{i\ell k} &\geq \alpha_{i\ell,k-1}
    - \rho_{i\ell}(\alpha_{i\ell,k-1} - \alpha^{\mathrm{last}}_{i\ell,k-1})
    - \bar{\alpha}_{i\ell}(1 - x^m_{i\ell k})
\end{align}

\textbf{Replacement} encoded as in ARD1 case.

\subsubsection{Reliability Constraint (Gamma, precomputed linear bound)}

Since $\Pr(D_{i\ell k} > \tau) = Q(\alpha_{i\ell k},\,\tau/\beta_{i\ell}) \leq \varepsilon$
is monotone increasing in $\alpha_{i\ell k}$, it reduces to:
\begin{equation}
  \alpha_{i\ell k} \leq \hat{\alpha}_{i\ell} \qquad \forall\, i,\ell,k
\end{equation}
where $\hat{\alpha}_{i\ell}$ is the unique solution to
$Q(\hat{\alpha}_{i\ell},\,\tau/\beta_{i\ell}) = \varepsilon$,
computed \textbf{offline} before the Gurobi model is built.
Note $\hat{\alpha}_{i\ell} \leq \bar{\alpha}_{i\ell}$ always, so the reliability
bound is strictly tighter than the big-M.

%% --- CHANGE DOM-016: Q convention and scipy mapping ---
\emph{Convention.} $Q(a,x) = 1 - P(a,x)$ is the \textbf{regularised upper} incomplete
gamma function, where $P(a,x) = \gamma(a,x)/\Gamma(a)$ is the regularised lower
incomplete gamma. In \texttt{scipy}: $Q(a,x)$ = \texttt{scipy.special.gammaincc(a, x)}
and $P(a,x)$ = \texttt{scipy.special.gammainc(a, x)}.

%% --- CHANGE CS-011: brentq root-find needed, not gammaincinv ---
\emph{Implementation.} Computing $\hat\alpha$ requires solving
$Q(\hat\alpha, \tau/\beta) = \varepsilon$ for $\hat\alpha$. This is a root-finding
problem in $\hat\alpha$, \textbf{not} a direct call to \texttt{gammaincinv} (which
solves for $x$, not~$a$). Use \texttt{scipy.optimize.brentq} with the objective
$g(\hat\alpha) = Q(\hat\alpha, \tau/\beta) - \varepsilon$.

\emph{Note.} If the root-finding yields $\hat\alpha > \tau/\beta$, the reliability
target is unachievable for this component given its degradation rate and threshold.
The implementer should raise a diagnostic warning.

\subsubsection{Loop Constraint (Gamma, hard + soft)}

%% --- CHANGE MATH-010: inline Gamma FSD proposition statement ---
By the FSD ordering result for Gamma distributions (Proposition~\texttt{prop:gamma\_fsd}
of the reference document): for $D_1 \sim \mathrm{Gamma}(\alpha_1, \beta)$ and
$D_2 \sim \mathrm{Gamma}(\alpha_2, \beta)$ with common scale $\beta$,
$D_1 \leq_{\mathrm{FSD}} D_2$ if and only if $\alpha_1 \leq \alpha_2$.
Therefore, FSD at fixed $\beta$ reduces to a shape condition:
\begin{equation}
  \alpha_{i\ell,2H} \leq \alpha_{i\ell,H}
  \qquad \forall\,(i,\ell)\text{ s.t. model} = \mathrm{Gamma}
\end{equation}
Loop penalty contribution:
\begin{equation}
  \mathcal{L}^{\mathrm{Gamma}}_{\mathrm{loop}} =
    \sum_{(i,\ell)\,:\,\mathrm{Gamma}} \bigl(\alpha_{i\ell H} - \alpha_{i\ell,2H}\bigr)
\end{equation}

% ----------------------------------------------------------------
%% --- CHANGE RF-023: new rainflow state-dynamics subsection ---
\subsection{State Dynamics --- Rainflow Model}
\label{sec:dyn_rainflow}
% ----------------------------------------------------------------

For components $(i,\ell)$ with \texttt{model} $=$ \texttt{rainflow}.
\textbf{Only \texttt{ARD1} is supported} (consistency check must enforce this):
as for the Gaussian model, an ARA1 extension acting on the variance state has no
established semantics: ARA1 is defined in the maintenance literature through the
virtual age, not through explicit variance tracking, so no principled
post-repair variance rule is available (see the Open Questions).
Both $\mu_{i\ell k}$ and $v_{i\ell k}$ are tracked.

The rainflow model is \textbf{distribution-free}: the per-mission
damage-increment law is measured empirically (rainflow counting of recorded load
histories) and enters the model only through its per-mission mean
$\mu^{\mathrm{inc}}_{ij\ell k}$, variance $v^{\mathrm{inc}}_{ij\ell k}$, and (for
the Bernstein certificate) support upper bound $b^{\mathrm{inc}}_{ij\ell k}$; no
parametric family is assumed. The reliability requirement
$\Pr(D_{i\ell k} > \tau) \leq \varepsilon$ is certified by a distribution-free
tail bound selected per component by the \texttt{tail\_bound} input
(\texttt{"cantelli"}, the default, or
\texttt{"bernstein"}).\footnote{The research grounding for this model---the
selection of tail bounds, the curvature analysis of the resulting constraints,
and the loop-on-moments argument---is developed in the project scoping
document, subsection ``Conservative surrogates for the empirical
damage-increment distribution''.}

\textbf{Modelling assumptions (normative).}
\begin{itemize}[leftmargin=1.5em]
  \item \textbf{A1 (independence).} All mission increments are mutually
        independent across days, independent of the initial damage state
        $D^{0}_{i\ell}$ and of every replacement draw; each replacement installs
        a fresh draw of $D^{\mathrm{new}}_{i\ell}$, independent of everything
        else. Without A1 the variance recursion below and both tail certificates
        are invalid.
  \item \textbf{A2 (bounded supports; Bernstein only).}
        $D^{0}_{i\ell} \in [0, b^{0}_{i\ell}]$ and
        $D^{\mathrm{new}}_{i\ell} \in [0, b^{\mathrm{new}}_{i\ell}]$ almost
        surely, with known finite bounds. Unbounded reset laws (e.g.\ a Gaussian
        $D^{\mathrm{new}}$) are inadmissible under the Bernstein certificate;
        the Cantelli certificate remains valid for them, needing only finite
        variance.
  \item \textbf{A3 (certificate coverage and periodicity).} The reliability
        certificate is imposed at \emph{every} step $k = 1,\ldots,2H$ on the
        tracked moments, and the per-mission increment parameters are
        $H$-periodic in $k$ (the $k \bmod H$ wrap of
        Section~\ref{sec:input_increments}). This
        periodicity is load-bearing, not an implementation convenience: it is
        what makes every future repetition of the second half execute the same
        moment map, which the loop constraint below requires.
\end{itemize}

\textbf{Notation.}
$\delta\mu_{i\ell k} = \sum_{j=1}^{M} x_{ijk}\,\mu^{\mathrm{inc}}_{ij\ell k}$,
$\delta v_{i\ell k} = \sum_{j=1}^{M} x_{ijk}\,v^{\mathrm{inc}}_{ij\ell k}$
(mission-induced increments; periodic extension for $k>H$), exactly as for the
Gaussian model. $\kappa_C = \sqrt{(1-\varepsilon)/\varepsilon}$ (Cantelli);
$\kappa = \ln(1/\varepsilon)$, $b_{i\ell}$, and
$c_{i\ell} = \kappa\, b_{i\ell}/3$ (Bernstein; see the reliability subsection).
Big-M values: $M_\mu = \tau$ (valid because the reliability constraint enforces
$\mu_{i\ell k} \leq \tau$, indeed $\mu_{i\ell k} \leq \tau - c_{i\ell}$ under
Bernstein) and $M_v = V_{\max}^{\mathrm{RF}}$ with
\begin{equation}
  V_{\max}^{\mathrm{RF}} =
  \begin{cases}
    \tau^2\,\varepsilon/(1-\varepsilon)
      & \texttt{tail\_bound="cantelli"} \\[2pt]
    \tau\,(\tau - 2\kappa\, b_{i\ell}/3)/(2\kappa)
      & \texttt{tail\_bound="bernstein"}
  \end{cases}
  \label{eq:rf_vmax}
\end{equation}
Each value is the largest variance the respective exact reliability constraint
admits, obtained at $\mu_{i\ell k} = 0$: for Cantelli,
$\kappa_C^2\, v \leq \tau^2$ gives $v \leq \tau^2/\kappa_C^2 =
\tau^2\varepsilon/(1-\varepsilon)$; for Bernstein,
$2\kappa\, v \leq (\tau - c_{i\ell})^2 - c_{i\ell}^2 = \tau(\tau - 2c_{i\ell})$
gives $v \leq \tau(\tau - 2\kappa\, b_{i\ell}/3)/(2\kappa)$.
The consistency check $\tau > 2\kappa\, b_{i\ell}/3$
(Section~\ref{sec:consistency_checks}) keeps the
Bernstein value positive.

\subsubsection{ARD1 Rainflow}

The moment dynamics are identical to the ARD1 Gaussian dynamics of
Section~\ref{sec:dyn_gaussian} (with $V_{\max}$ replaced by
$V_{\max}^{\mathrm{RF}}$): they are pure moment-propagation bookkeeping and carry
no distributional assumption.

\textbf{Mean dynamics --- exact form:}
\begin{equation}
  \mu_{i\ell k} = \begin{cases}
    \mu_{i\ell,k-1} + \delta\mu_{i\ell k} & \text{if } x^m_{i\ell k} = 0 \\
    (1-\rho_{i\ell})\,\mu_{i\ell,k-1} & \text{if } x^m_{i\ell k} = 1
  \end{cases}
\end{equation}

\textbf{Mean dynamics --- linearised (big-M):}
\begin{align}
  \mu_{i\ell k} &\geq \mu_{i\ell,k-1} + \delta\mu_{i\ell k} - \tau\,x^m_{i\ell k}
    - \tau\,x^r_{i\ell k}
    \label{eq:rf_ard_mu1} \\
  \mu_{i\ell k} &\geq (1-\rho_{i\ell})\,\mu_{i\ell,k-1} - \tau(1-x^m_{i\ell k})
    \label{eq:rf_ard_mu2}
\end{align}

\textbf{Variance dynamics --- exact form:}
\begin{equation}
  v_{i\ell k} = \begin{cases}
    v_{i\ell,k-1} + \delta v_{i\ell k} & \text{if } x^m_{i\ell k} = 0 \\
    (1-\rho_{i\ell})^2\,v_{i\ell,k-1} & \text{if } x^m_{i\ell k} = 1
  \end{cases}
\end{equation}

\textbf{Variance dynamics --- linearised (big-M):}
\begin{align}
  v_{i\ell k} &\geq v_{i\ell,k-1} + \delta v_{i\ell k}
    - V_{\max}^{\mathrm{RF}}\,x^m_{i\ell k}
    - V_{\max}^{\mathrm{RF}}\,x^r_{i\ell k} \\
  v_{i\ell k} &\geq (1-\rho_{i\ell})^2\,v_{i\ell,k-1}
    - V_{\max}^{\mathrm{RF}}(1-x^m_{i\ell k})
\end{align}

\textbf{Replacement --- exact form:}
\begin{equation}
  (\mu_{i\ell k},\, v_{i\ell k}) = (\mu^{\mathrm{new}}_{i\ell},\, v^{\mathrm{new}}_{i\ell})
  \quad \text{if } x^r_{i\ell k} = 1
  \label{eq:rf_rep}
\end{equation}

\textbf{Replacement --- linearised (big-M):}
\begin{align}
  \mu_{i\ell k} &\leq \mu^{\mathrm{new}}_{i\ell} + \tau(1-x^r_{i\ell k}), \quad
  \mu_{i\ell k} \geq \mu^{\mathrm{new}}_{i\ell} - \tau(1-x^r_{i\ell k}) \\
  v_{i\ell k} &\leq v^{\mathrm{new}}_{i\ell} + V_{\max}^{\mathrm{RF}}(1-x^r_{i\ell k}), \quad
  v_{i\ell k} \geq v^{\mathrm{new}}_{i\ell} - V_{\max}^{\mathrm{RF}}(1-x^r_{i\ell k})
\end{align}

\emph{Remark (moment exactness).}
Under assumption A1 the tracked pair $(\mu_{i\ell k}, v_{i\ell k})$ equals the
\emph{true} mean and variance of the accumulated damage state at every step,
whatever the underlying increment laws: a mission day adds the increment's mean
and variance (independence makes variances additive); ARD1 repair applies
$D^+ = (1-\rho)\,D^-$, hence $\mu^+ = (1-\rho)\,\mu^-$ and
$v^+ = (1-\rho)^2\,v^-$ exactly; replacement resets both moments to those of the
fresh draw $(\mu^{\mathrm{new}}_{i\ell}, v^{\mathrm{new}}_{i\ell})$; idle days
change nothing. The recursions are exact---not approximations---and the big-M
linearisation is exact at binary points. The tightness argument of
Section~\ref{sec:dyn_gaussian} (three cost mechanisms for $\mu$; reliability
pressure for $v$) applies unchanged.

\subsubsection{Reliability Constraint (Rainflow)}

The certificate is selected per component by \texttt{tail\_bound} and imposed at
every step $k = 1,\ldots,2H$ on the tracked moments
$(\mu_{i\ell k}, v_{i\ell k})$. Because the tracked moments are exact (remark
above), any point satisfying the constraint carries a valid distribution-free
certificate $\Pr(D_{i\ell k} > \tau) \leq \varepsilon$.

\textbf{Cantelli bound (\texttt{tail\_bound="cantelli"}).}
For any law with mean $\mu$, variance $v$, and $\tau > \mu$, Cantelli's
inequality gives $\Pr(D > \tau) \leq v/(v + (\tau - \mu)^2)$; this bound is at
most $\varepsilon$ if and only if
\begin{equation}
  \mu_{i\ell k} + \kappa_C\,\sqrt{v_{i\ell k}} \leq \tau
  \qquad \forall\, i,\ell,k,
  \qquad \kappa_C = \sqrt{(1-\varepsilon)/\varepsilon}
  \label{eq:rf_cantelli}
\end{equation}
(rearranging $v \leq \varepsilon\,(v + (\tau-\mu)^2)$ gives
$(1-\varepsilon)\,v \leq \varepsilon\,(\tau-\mu)^2$, i.e.\
$\tau - \mu \geq \kappa_C\sqrt{v}$). Only the two tracked moments are required;
no independence or boundedness assumption enters the bound itself (assumption A1
is still required for the tracked $v_{i\ell k}$ to equal the true variance).

\emph{Exact form (\texttt{formulation="exact"}):}
\begin{equation}
  \kappa_C^2\, v_{i\ell k} \leq (\tau - \mu_{i\ell k})^2,
  \qquad \mu_{i\ell k} \leq \tau
  \qquad \forall\, i,\ell,k
  \label{eq:rf_cantelli_exact}
\end{equation}
The linear row keeps $\tau - \mu_{i\ell k} \geq 0$, so the squared form is
equivalent to Eq.~\eqref{eq:rf_cantelli}. This is structurally identical to the
Gaussian exact form with $(\Phi^{-1}(1-\varepsilon))^2$ replaced by
$\kappa_C^2 = (1-\varepsilon)/\varepsilon$, and, like it, is a \textbf{nonconvex
quadratic} constraint, \emph{not} a rotated SOC (the $\mu_{i\ell k}^2$ term
enters with the wrong sign for convexity; see the corrected discussion in
Section~\ref{sec:dyn_gaussian}). Gurobi handles it via its nonconvex quadratic
support (\texttt{NonConvex=2} where not the default); the certificate is exact.

\emph{LP surrogate (\texttt{formulation="lp"}):} expanding
$(\tau - \mu_{i\ell k})^2$ and dropping the non-negative term $\mu_{i\ell k}^2$:
\begin{equation}
  \kappa_C^2\, v_{i\ell k} + 2\tau\,\mu_{i\ell k} \leq \tau^2
  \qquad \forall\, i,\ell,k
  \label{eq:rf_cantelli_lp}
\end{equation}
This is an \textbf{inner approximation}, derived exactly as the Gaussian LP form
(Eq.~\eqref{eq:g_lp_reliability}): any point satisfying
Eq.~\eqref{eq:rf_cantelli_lp} satisfies Eq.~\eqref{eq:rf_cantelli_exact}, since
$(\tau - \mu_{i\ell k})^2 = \tau^2 - 2\tau\,\mu_{i\ell k} + \mu_{i\ell k}^2
\geq \tau^2 - 2\tau\,\mu_{i\ell k}$. The approximation error is exactly
$\mu_{i\ell k}^2$, small when $\mu_{i\ell k} \ll \tau$. The conservatism remark
of the Gaussian LP form applies verbatim: the surrogate may exclude
exact-feasible points and can only increase the objective; it never invalidates
the certificate.

\textbf{Bernstein bound (\texttt{tail\_bound="bernstein"}).}
Write $\kappa = \ln(1/\varepsilon)$ and define the per-component support
constant
\begin{equation}
  b_{i\ell} = \max\Bigl\{\, b^{0}_{i\ell},\; b^{\mathrm{new}}_{i\ell},\;
    \max_{j,k}\, b^{\mathrm{inc}}_{ij\ell k} \,\Bigr\},
  \qquad
  c_{i\ell} = \frac{\kappa\, b_{i\ell}}{3}.
  \label{eq:rf_bsupport}
\end{equation}
$b_{i\ell}$ is a \textbf{precomputed constant}, not a decision variable: under
ARD1, between replacements the damage state unrolls as a weighted sum
$w_0\, D^{\mathrm{reset}} + \sum_t w_t\, \Delta D_t$ whose repair weights
$w \in [0,1]$ only \emph{shrink} each summand's support, so every summand is
bounded above by $b_{i\ell}$ regardless of the schedule. Taking the maximum over
all admissible summands---initial state, replacement draws, and every mission
increment---therefore yields a valid, schedule-independent support bound
(conservative to the extent that the maximum ranges over missions not actually
served).

For a sum of mutually independent summands each bounded above by $b_{i\ell}$
(assumptions A1--A2)---each summand is also non-negative, so its centered
version satisfies $X - \mathbb{E}[X] \leq X \leq b_{i\ell}$ a fortiori, as
Bernstein's inequality requires---with mean $\mu$ and variance $v$, Bernstein's
inequality
gives $\Pr(D > \tau) \leq
\exp\bigl(-a^2/(2(v + b_{i\ell}\,a/3))\bigr)$ with $a = \tau - \mu > 0$.
Requiring the right-hand side to be at most $\varepsilon$ is equivalent to
$a^2 - 2c_{i\ell}\,a - 2\kappa\, v \geq 0$, i.e.\ (larger root of the quadratic)
$a \geq c_{i\ell} + \sqrt{c_{i\ell}^2 + 2\kappa\, v}$, yielding the chance-bound
form
\begin{equation}
  \mu_{i\ell k} + c_{i\ell}
  + \sqrt{c_{i\ell}^2 + 2\kappa\, v_{i\ell k}} \;\leq\; \tau
  \qquad \forall\, i,\ell,k,
  \label{eq:rf_bernstein}
\end{equation}
which is sufficient for $\Pr(D_{i\ell k} > \tau) \leq \varepsilon$ (the
precondition $a > 0$ is self-enforced, since the left-hand side exceeds
$\mu_{i\ell k}$ whenever $b_{i\ell} > 0$ or $v_{i\ell k} > 0$).

\emph{Exact form (\texttt{formulation="exact"}):}
\begin{equation}
  (\tau - c_{i\ell} - \mu_{i\ell k})^2 \geq c_{i\ell}^2 + 2\kappa\, v_{i\ell k},
  \qquad \mu_{i\ell k} \leq \tau - c_{i\ell}
  \qquad \forall\, i,\ell,k
  \label{eq:rf_bernstein_exact}
\end{equation}
The linear row makes the bracket non-negative, so the squared form is equivalent
to Eq.~\eqref{eq:rf_bernstein}. As for Cantelli, this is a nonconvex quadratic
constraint; the same Gurobi note applies.

\emph{LP surrogate (\texttt{formulation="lp"}):} expanding
$(\tau - c_{i\ell} - \mu_{i\ell k})^2$ and dropping $\mu_{i\ell k}^2 \geq 0$,
then using $(\tau - c_{i\ell})^2 - c_{i\ell}^2 = \tau\,(\tau - 2c_{i\ell})$:
\begin{equation}
  2(\tau - c_{i\ell})\,\mu_{i\ell k} + 2\kappa\, v_{i\ell k}
  \;\leq\; \tau\,(\tau - 2c_{i\ell})
  \qquad \forall\, i,\ell,k
  \label{eq:rf_bernstein_lp}
\end{equation}
This is the tangent at $\mu_{i\ell k} = 0$ of the convex parabola (in
$\mu_{i\ell k}$) defined by Eq.~\eqref{eq:rf_bernstein_exact}; the derivation
and the conservatism are exactly parallel to the Cantelli and Gaussian LP forms,
with approximation error again exactly $\mu_{i\ell k}^2$. The feasibility
precondition $\tau > 2\kappa\, b_{i\ell}/3$
(Section~\ref{sec:consistency_checks}) keeps the right-hand
side positive.

\emph{Remark (tighter piecewise-linear surrogates).}
The single-tangent surrogates above can in principle be tightened to the
pointwise minimum of several tangents, which is still a global over-estimator of
the square root and hence still conservative. This, however, requires encoding
the piecewise-linear function as an \emph{equality} (SOS2 $\lambda$-formulation
or Gurobi's \texttt{addGenConstrPWL}): a mere conjunction of tangent
inequalities on an auxiliary variable is either over-conservative (imposing
$y \geq$ every tangent realises their pointwise \emph{maximum}, worse than any
single tangent) or unsafe (imposing $y \leq$ every tangent lets the auxiliary
under-estimate the square root, breaking conservatism). Deferred; see the Open
Questions.

\textbf{User-defined variance upper bound (optional):}
If $V_{\max}^{\mathrm{user}}$ is specified for rainflow component $(i,\ell)$:
\begin{equation}
  v_{i\ell k} \leq V_{\max,i\ell}^{\mathrm{user}} \qquad \forall\, k
\end{equation}

\subsubsection{Loop Constraint (Rainflow, hard + soft)}

Hard constraints:
\begin{equation}
  \mu_{i\ell,2H} \leq \mu_{i\ell,H}, \qquad
  v_{i\ell,2H} \leq v_{i\ell,H}
  \qquad \forall\,(i,\ell)\text{ s.t. model} = \mathrm{Rainflow}
  \label{eq:rf_loop}
\end{equation}
Loop penalty contribution:
\begin{equation}
  \mathcal{L}^{\mathrm{Rainflow}}_{\mathrm{loop}} =
    \sum_{(i,\ell)\,:\,\mathrm{Rainflow}}
    \bigl[(\mu_{i\ell H} - \mu_{i\ell,2H}) + (v_{i\ell H} - v_{i\ell,2H})\bigr]
\end{equation}

\emph{Justification (moment domination chains; the scalar bound does not).}
The certified failure-probability bound $U(\mu, v)$---Cantelli
$v/(v + (\tau - \mu)^2)$ or Bernstein
$\exp\bigl(-(\tau-\mu)^2/(2(v + b_{i\ell}(\tau-\mu)/3))\bigr)$---is increasing
in $v$ and, for $\mu < \tau$, increasing in $\mu$. Every per-day moment
operation is componentwise non-decreasing in $(\mu, v)$: a mission increment
adds non-negative constants; ARD1 repair multiplies by the non-negative factors
$(1-\rho)$ and $(1-\rho)^2$; replacement maps to constants; idle days are the
identity. Compositions of componentwise non-decreasing maps are componentwise
non-decreasing, so the second-half period map $\Phi$ (days $H{+}1,\ldots,2H$)
and every prefix $\Phi_r$ of it are monotone. Eq.~\eqref{eq:rf_loop} states
$\Phi(\mu_H, v_H) \leq (\mu_H, v_H)$ componentwise; under indefinite
concatenation of the second half, induction gives
$(\mu_{(m+1)H}, v_{(m+1)H}) = \Phi(\mu_{mH}, v_{mH}) \leq
\Phi(\mu_{(m-1)H}, v_{(m-1)H}) = (\mu_{mH}, v_{mH})$ for every $m \geq 2$,
the case $m = 1$ being precisely Eq.~\eqref{eq:rf_loop}; hence
$(\mu_{mH}, v_{mH}) \leq (\mu_H, v_H)$ for all $m \geq 1$, and
for intermediate days
$(\mu_{mH+r}, v_{mH+r}) = \Phi_r(\mu_{mH}, v_{mH}) \leq
\Phi_r(\mu_H, v_H) = (\mu_{H+r}, v_{H+r})$. Since the certificate is imposed at
every nominal step $k = 1,\ldots,2H$ (assumption A3), every step of every
future repetition carries a certified bound
$U_{mH+r} \leq U_{H+r} \leq \varepsilon$. The $H$-periodicity of the increment
parameters (A3) is what makes each repetition execute the same map $\Phi$.

\emph{Warning (do not loop on the scalar bound).}
Imposing the single comparison $U_{2H} \leq U_H$ in place of the two moment
inequalities would \emph{not} chain: $U$ is a many-to-one function of
$(\mu, v)$ whose level sets are not invariant under the dynamics, so
$U_{2H} \leq U_H$ can hold with a larger mean offset by a smaller variance, and
one further period can produce $U_{3H} > U_H$.\footnote{A concrete instance:
$\tau = 10$, $\rho = 0.5$, repeated block $=$ one repair day followed by mission
days with block totals $(\sum\delta\mu, \sum\delta v) = (3, 3.99)$; from
$(\mu_H, v_H) = (4, 9)$ the Cantelli bound gives $U_H = 0.2000$ and
$U_{2H} = 0.1997 \leq U_H$, yet $U_{3H} = 0.2151 > U_H$, with
$U_{mH} \to 0.2495$. An analogous Bernstein instance ($b_{i\ell} = 3$, block
variance $3.15$) also violates at $m = 3$. The moment condition
Eq.~\eqref{eq:rf_loop} rejects both instances, since
$\mu_{2H} = 5 > 4 = \mu_H$.} Eq.~\eqref{eq:rf_loop} is strictly stronger, and
that strengthening is exactly what makes the induction valid. Note that
Eq.~\eqref{eq:rf_loop} is a statement of \emph{moment domination} yielding a
non-increasing certified failure-probability bound; it is not a
stochastic-order (e.g.\ FSD) claim about the underlying damage laws.

% ----------------------------------------------------------------
\subsection{Damage Regularisation Variable}
\label{sec:objective}
% ----------------------------------------------------------------

$u$ is a single non-negative scalar capturing the worst per-train, per-step
aggregate damage across the entire horizon. Its definition depends on
\texttt{penalty\_type}.

\subsubsection{\texttt{inf\_norm} (default, MILP)}

\textbf{Exact form:}
\begin{equation}
  u = \max_{i,k} \sum_{\ell=1}^{L} \mu_{i\ell k}
\end{equation}

\textbf{Linearised form:}
\begin{equation}
  u \geq \sum_{\ell=1}^{L} \mu_{i\ell k} \qquad \forall\, i,k
\end{equation}
This is linear in the decision variables; the problem remains an MILP.
For Gamma components, $\mu_{i\ell k} = \alpha_{i\ell k}\cdot\beta_{i\ell}$
is substituted directly.

%% --- CHANGE RF-024: formulation token socp->exact; penalty is genuinely convex SOC ---
\subsubsection{\texttt{quadratic} (convex rotated SOC; requires \texttt{formulation="exact"})}

\textbf{Exact form:}
\begin{equation}
  u = \max_{i,k} \sum_{\ell=1}^{L} \mu_{i\ell k}^2
\end{equation}

\textbf{Linearised form:}
\begin{equation}
  u \geq \sum_{\ell=1}^{L} \mu_{i\ell k}^2 \qquad \forall\, i,k
\end{equation}
Each constraint $u \geq \sum_\ell \mu_{i\ell k}^2$ is a rotated SOC constraint
($u \cdot 1 \geq \|\boldsymbol{\mu}_{ik}\|_2^2$ with $x=u$, $y=1$,
$\|z\|_2^2 = \sum_\ell \mu_{i\ell k}^2$),
natively supported by Gurobi. Unlike the exact reliability constraints of the
Gaussian, Wiener, and rainflow models---which are \emph{nonconvex}
quadratics (Section~\ref{sec:dyn_gaussian})---this penalty constraint is
genuinely convex: the quadratic terms appear on the smaller side of the
inequality. Under \texttt{formulation="exact"} the overall problem is already a
nonconvex MIQCP \emph{when Gaussian, Wiener, or rainflow components are
present}, so in that case the convex quadratic penalty adds no new problem
class (with only IG/Gamma components it upgrades the otherwise-linear model to
a convex MIQCP); it is incompatible with \texttt{formulation="lp"} because it
is not linear.

% ----------------------------------------------------------------
\subsection{Repair Cost Variable}
\label{sec:repair_proxy}
% ----------------------------------------------------------------

$z_{i\ell k}$ represents the degradation removed by imperfect repair on component
$\ell$ of train $i$ at step $k$, expressed in mean units. It is the product of the
repair efficiency, the pre-repair mean, and the repair binary:

%% --- CHANGE DEV-016: z is true degradation removed, not a proxy ---
\textbf{Exact (bilinear) form --- ARD1:}
\begin{equation}
  z_{i\ell k} = \rho_{i\ell}\,\mu^-_{i\ell,k}\,x^m_{i\ell k}
  \qquad \forall\, i,\ell,k \text{ (ARD1 components)}
\end{equation}

\textbf{Exact (bilinear) form --- ARA1:}
\begin{equation}
  z_{i\ell k} = \rho_{i\ell}\,(\mu^-_{i\ell,k} - \mu^{\mathrm{last}}_{i\ell,k})\,x^m_{i\ell k}
  \qquad \forall\, i,\ell,k \text{ (ARA1 components)}
\end{equation}

%% --- CHANGE MATH-012: Gamma alpha*beta substitution note ---
\emph{Gamma substitution.} For Gamma components, $\mu_{i\ell,k-1}$ in the McCormick
bounds should be read as $\alpha_{i\ell,k-1} \cdot \beta_{i\ell}$, since the Gamma
model tracks the shape parameter $\alpha$ rather than the mean directly.

\textbf{Linearised form (McCormick envelope, big-M $= \rho_{i\ell}\,\tau$; ARD1 shown):}
\begin{align}
  z_{i\ell k} &\leq \rho_{i\ell}\,\mu_{i\ell,k-1} \\
  z_{i\ell k} &\leq \rho_{i\ell}\,\tau\,x^m_{i\ell k} \\
  z_{i\ell k} &\geq \rho_{i\ell}\,\mu_{i\ell,k-1} - \rho_{i\ell}\,\tau\,(1-x^m_{i\ell k}) \\
  z_{i\ell k} &\geq 0
\end{align}
Note: $\rho_{i\ell}\,\tau$ is a valid big-M since $\mu_{i\ell,k-1} \leq \tau$ always.

% ----------------------------------------------------------------
\subsection{Loop Constraint Penalty}
\label{sec:loop}
% ----------------------------------------------------------------

The penalty $\mathcal{L}_{\mathrm{loop}}$ enforces sustainability over the two
half-horizons $[1,H]$ and $[H{+}1,2H]$.
Both a \textbf{hard constraint} and a \textbf{soft penalty} are applied for each model.

\subsubsection{Gaussian}
Hard constraints:
\begin{equation}
  \mu_{i\ell,2H} \leq \mu_{i\ell,H}, \qquad
  v_{i\ell,2H} \leq v_{i\ell,H} \qquad
  \forall\, (i,\ell)\text{ s.t. model} = \mathrm{Gaussian}
\end{equation}
Penalty:
\begin{equation}
  \mathcal{L}^{\mathrm{Gauss}}_{\mathrm{loop}} =
    \sum_{(i,\ell)\,:\,\mathrm{Gauss}}
    \bigl[(\mu_{i\ell H} - \mu_{i\ell 2H}) + (v_{i\ell H} - v_{i\ell 2H})\bigr]
\end{equation}

\subsubsection{Inverse Gaussian}
Hard constraint and penalty as defined in Section~\ref{sec:dyn_ig}.

\subsubsection{Wiener}
%% --- CHANGE RF-025: A-02 accepted review item applied (FSD wording corrected) ---
Hard constraints and penalty as defined in Section~\ref{sec:dyn_wiener}.
Both the mean and the variance state are bounded (implying a non-increasing
failure probability for Gaussian marginals with $\mu < \tau$; not an FSD claim);
variance control requires replacement actions within $[H+1,2H]$.

\subsubsection{Gamma}
Hard constraint and penalty as defined in Section~\ref{sec:dyn_gamma}.

%% --- CHANGE RF-026: rainflow loop cross-reference and total penalty term ---
\subsubsection{Rainflow}
Hard constraints and penalty as defined in Section~\ref{sec:dyn_rainflow}.
Both tracked moments are bounded, which guarantees a non-increasing
\emph{certified} failure-probability bound under indefinite schedule
concatenation (moment domination; not a distributional-order claim).

The total loop penalty is:
\begin{equation}
  \mathcal{L}_{\mathrm{loop}} =
    \mathcal{L}^{\mathrm{Gauss}}_{\mathrm{loop}}
    + \mathcal{L}^{\mathrm{IG}}_{\mathrm{loop}}
    + \mathcal{L}^{\mathrm{Wiener}}_{\mathrm{loop}}
    + \mathcal{L}^{\mathrm{Gamma}}_{\mathrm{loop}}
    + \mathcal{L}^{\mathrm{Rainflow}}_{\mathrm{loop}}
\end{equation}

%% --- CHANGE MATH-016: loop penalty non-negativity remark ---
\emph{Remark (non-negativity of $\mathcal{L}_{\mathrm{loop}}$).}
Each term $(\mu_{i\ell H} - \mu_{i\ell 2H})$ (and analogously for $v$ and $\alpha$)
in the loop penalty is non-negative at any feasible point, because the hard loop
constraints enforce $\mu_{i\ell,2H} \leq \mu_{i\ell,H}$ (and likewise for $v$, $\alpha$).
The soft penalty therefore acts as a tie-breaker among feasible schedules, favouring
those with the largest sustainability margin.

% ================================================================
\section{Full Optimization Problem (\texttt{inf\_norm}, linearised)}
\label{sec:full_problem}
% ================================================================
\thispagestyle{fancy}

%% --- CHANGE CS-002: constraint count table ---
\begin{table}[htbp]
\centering
\caption{Approximate constraint count per $(i,\ell,k)$ cell by model and maintenance type.
Multiply by $F \times L \times 2H$ for the full model.
Global constraints (assignment, mission coverage) add $O(F \times M \times 2H)$.}
\begin{tabular}{lccc}
  \toprule
  Model--Maintenance & Binary vars & Continuous vars & Constraints \\
  \midrule
  Gaussian--ARD1     & $M+1$ & $3 + M$ & $\sim 12$ \\
  Wiener--ARD1       & $M+1$ & $3 + M$ & $\sim 10$ \\
  Wiener--ARA1       & $M+1$ & $4 + M$ & $\sim 14$ \\
  IG--ARD1           & $M+1$ & $2 + M$ & $\sim 10$ \\
  Gamma--ARD1        & $M+1$ & $2 + M$ & $\sim 10$ \\
  Gamma--ARA1        & $M+1$ & $3 + M$ & $\sim 14$ \\
  %% --- CHANGE RF-027: rainflow constraint-count row ---
  Rainflow--ARD1     & $M+1$ & $3 + M$ & $\sim 12$ \\
  \bottomrule
\end{tabular}
\label{tab:constraint_count}
\end{table}

\begin{small}
\allowdisplaybreaks

\textbf{Notation used below.}
$\delta\mu_{i\ell k} = \sum_{j=1}^{M} x_{ijk}\,\mu^{\mathrm{inc}}_{ij\ell k}$;
%% --- CHANGE RF-028: rainflow notation in full problem ---
$\delta v_{i\ell k} = \sum_{j=1}^{M} x_{ijk}\,v^{\mathrm{inc}}_{ij\ell k}$ (Gaussian, Rainflow);
$\delta v^W_{i\ell k} = \sigma^2_{i\ell}\sum_{j=1}^{M} x_{ijk}$ (Wiener);
$\delta\alpha_{i\ell k} = \sum_{j=1}^{M} x_{ijk}\,\alpha^{\mathrm{inc}}_{ij\ell k}$ (Gamma);
$V_{\max} = \tau^2/(\Phi^{-1}(1-\varepsilon))^2$;
$\bar{\alpha}_{i\ell} = \tau/\beta_{i\ell}$;
$\kappa_C = \sqrt{(1-\varepsilon)/\varepsilon}$,
$\kappa = \ln(1/\varepsilon)$,
$b_{i\ell} = \max\{b^{0}_{i\ell},\, b^{\mathrm{new}}_{i\ell},\,
  \max_{j,k} b^{\mathrm{inc}}_{ij\ell k}\}$,
$c_{i\ell} = \kappa\, b_{i\ell}/3$ (Rainflow);
$V_{\max}^{\mathrm{RF}} = \tau^2\varepsilon/(1-\varepsilon)$ (Cantelli) or
$\tau(\tau - 2c_{i\ell})/(2\kappa)$ (Bernstein).
Superscripts G, IG, W, $\Gamma$, RF restrict constraints to the respective model components.
For rainflow components, the Cantelli rows apply when
\texttt{tail\_bound="cantelli"} and the Bernstein rows when
\texttt{tail\_bound="bernstein"}; exact rows apply under
\texttt{formulation="exact"}, LP rows under \texttt{formulation="lp"}.

\vspace{0.5em}
\noindent\textbf{Objective}
\begin{align}
\min\quad
  & C_M \sum_{k,i} x_{i0k}
  + \sum_{k,i,\ell} C_{R,i\ell}\,z_{i\ell k}
  + \sum_{k,i,\ell} C_{\mathrm{rep},i\ell}\,x^r_{i\ell k}
  + C_D\,u
  + C_P\,\mathcal{L}_{\mathrm{loop}} \tag{OBJ}
\end{align}

\vspace{0.3em}
\noindent\textbf{Assignment}
\begin{alignat}{2}
  &\textstyle\sum_{j=0}^{M} x_{ijk} \leq 1
  &&\quad\forall\,i,k \tag{one activity per train} \\
  &\textstyle\sum_{i=1}^{F} x_{ijk} = 1
  &&\quad\forall\,j\geq 1,\,k \tag{mission coverage} \\
  &x^m_{i\ell k} \leq x_{i0k},\quad x^r_{i\ell k} \leq x_{i0k}
  &&\quad\forall\,i,\ell,k \tag{maintenance on maint.\ day} \\
  &x^m_{i\ell k} + x^r_{i\ell k} \leq 1
  &&\quad\forall\,i,\ell,k \tag{no simultaneous repair+replace}
\end{alignat}

%% --- CHANGE MATH-019: add -tau*x^r / -V_max*x^r / -bar_alpha*x^r throughout full problem ---
\vspace{0.3em}
%% --- CHANGE MATH-019+DOM-009: ARD1 only for Gaussian, per-type annotations ---
\noindent\textbf{Gaussian dynamics} ($\forall\,(i,\ell)$ s.t.\ model$=$G; ARD1 only)
\begin{alignat}{2}
  &\mu_{i\ell k} \geq \mu_{i\ell,k-1} + \delta\mu_{i\ell k} - \tau\,x^m_{i\ell k}
    - \tau\,x^r_{i\ell k}
  &&\quad\forall\,k \tag{G: mean accumulation} \\
  &\mu_{i\ell k} \geq (1-\rho_{i\ell})\mu_{i\ell,k-1} - \tau(1-x^m_{i\ell k})
  &&\quad\forall\,k \tag{G: ARD1 mean repair} \\
  &\mu_{i\ell k} \leq \mu^{\mathrm{new}}_{i\ell} + \tau(1-x^r_{i\ell k}),\quad
   \mu_{i\ell k} \geq \mu^{\mathrm{new}}_{i\ell} - \tau(1-x^r_{i\ell k})
  &&\quad\forall\,k \tag{G: mean replacement} \\
  &v_{i\ell k} \geq v_{i\ell,k-1} + \delta v_{i\ell k} - V_{\max}\,x^m_{i\ell k}
    - V_{\max}\,x^r_{i\ell k}
  &&\quad\forall\,k \tag{G: var.\ accumulation} \\
  &v_{i\ell k} \geq (1-\rho_{i\ell})^2 v_{i\ell,k-1} - V_{\max}(1-x^m_{i\ell k})
  &&\quad\forall\,k \tag{G: ARD1 var.\ repair} \\
  &v_{i\ell k} \leq v^{\mathrm{new}}_{i\ell} + V_{\max}(1-x^r_{i\ell k}),\quad
   v_{i\ell k} \geq v^{\mathrm{new}}_{i\ell} - V_{\max}(1-x^r_{i\ell k})
  &&\quad\forall\,k \tag{G: var.\ replacement} \\
  %% --- CHANGE RF-029: reliability tag corrected (exact nonconvex quadratic, not SOC) ---
  &\mu_{i\ell k} + \Phi^{-1}(1-\varepsilon)\sqrt{v_{i\ell k}} \leq \tau
  &&\quad\forall\,k \tag{G: reliability (exact, nonconvex quad.)} \\
  &(\Phi^{-1}(1-\varepsilon))^2 v_{i\ell k} + 2\tau\,\mu_{i\ell k} \leq \tau^2
  &&\quad\forall\,k \tag{G: reliability (LP approx., inner)} \\
  &\mu_{i\ell,2H} \leq \mu_{i\ell,H},\quad v_{i\ell,2H} \leq v_{i\ell,H}
  &&\tag{G: loop constraint}
\end{alignat}

\vspace{0.3em}
\noindent\textbf{Inverse Gaussian dynamics} ($\forall\,(i,\ell)$ s.t.\ model$=$IG; ARD1 only)
\begin{alignat}{2}
  &\mu_{i\ell k} \geq \mu_{i\ell,k-1} + \delta\mu_{i\ell k} - \tau\,x^m_{i\ell k}
    - \tau\,x^r_{i\ell k}
  &&\quad\forall\,k \tag{IG: mean accumulation} \\
  &\mu_{i\ell k} \geq (1-\rho_{i\ell})\mu_{i\ell,k-1} - \tau(1-x^m_{i\ell k})
  &&\quad\forall\,k \tag{IG: ARD1 mean repair} \\
  &\mu_{i\ell k} \leq \mu^{\mathrm{new}}_{i\ell} + \tau(1-x^r_{i\ell k}),\quad
   \mu_{i\ell k} \geq \mu^{\mathrm{new}}_{i\ell} - \tau(1-x^r_{i\ell k})
  &&\quad\forall\,k \tag{IG: replacement} \\
  &\mu_{i\ell k} \leq \bar{\mu}_{i\ell}
  &&\quad\forall\,k \tag{IG: reliability (precomputed)} \\
  &\mu_{i\ell,2H} \leq \mu_{i\ell,H}
  &&\tag{IG: loop constraint}
\end{alignat}

\vspace{0.3em}
\noindent\textbf{Wiener dynamics} ($\forall\,(i,\ell)$ s.t.\ model$=$W)
\begin{alignat}{2}
  &\mu_{i\ell k} \geq \mu_{i\ell,k-1} + \delta\mu_{i\ell k} - \tau\,x^m_{i\ell k}
    - \tau\,x^r_{i\ell k}
  &&\quad\forall\,k \tag{W: mean accumulation} \\
  &\mu_{i\ell k} \geq (1-\rho_{i\ell})\mu_{i\ell,k-1} - \tau(1-x^m_{i\ell k})
  &&\quad\forall\,k \tag{W: ARD1 mean repair (ARD1); ARA1 uses $\mu^{\mathrm{last}}$} \\
  &\mu_{i\ell k} \leq \mu^{\mathrm{new}}_{i\ell} + \tau(1-x^r_{i\ell k}),\quad
   \mu_{i\ell k} \geq \mu^{\mathrm{new}}_{i\ell} - \tau(1-x^r_{i\ell k})
  &&\quad\forall\,k \tag{W: mean replacement} \\
  &\mu^{\mathrm{last}}_{i\ell k} \geq \mu_{i\ell,k-1} - \tau(1-x^m_{i\ell k}),\quad
   \mu^{\mathrm{last}}_{i\ell k} \leq \mu_{i\ell,k-1} + \tau(1-x^m_{i\ell k})
  &&\quad\forall\,k \tag{W: ARA1 $\mu^{\mathrm{last}}$ update (i)} \\
  &\mu^{\mathrm{last}}_{i\ell k} \geq \mu^{\mathrm{last}}_{i\ell,k-1} - \tau\,x^m_{i\ell k},\quad
   \mu^{\mathrm{last}}_{i\ell k} \leq \mu^{\mathrm{last}}_{i\ell,k-1} + \tau\,x^m_{i\ell k}
  &&\quad\forall\,k \tag{W: ARA1 $\mu^{\mathrm{last}}$ update (ii)} \\
  &v_{i\ell k} \geq v_{i\ell,k-1} + \delta v^W_{i\ell k} - V_{\max}\,x^r_{i\ell k}
  &&\quad\forall\,k \tag{W: var.\ accumulation (missions only)} \\
  &v_{i\ell k} \leq v^{\mathrm{new}}_{i\ell} + V_{\max}(1-x^r_{i\ell k}),\quad
   v_{i\ell k} \geq v^{\mathrm{new}}_{i\ell} - V_{\max}(1-x^r_{i\ell k})
  &&\quad\forall\,k \tag{W: var.\ replacement} \\
  %% --- CHANGE RF-030: reliability tag corrected (exact nonconvex quadratic, not SOC) ---
  &\mu_{i\ell k} + \Phi^{-1}(1-\varepsilon)\sqrt{v_{i\ell k}} \leq \tau
  &&\quad\forall\,k \tag{W: reliability (exact, nonconvex quad.)} \\
  &(\Phi^{-1}(1-\varepsilon))^2 v_{i\ell k} + 2\tau\,\mu_{i\ell k} \leq \tau^2
  &&\quad\forall\,k \tag{W: reliability (LP approx., inner)} \\
  &\mu_{i\ell,2H} \leq \mu_{i\ell,H},\quad v_{i\ell,2H} \leq v_{i\ell,H}
  &&\tag{W: loop constraint}
\end{alignat}

\vspace{0.3em}
\noindent\textbf{Gamma dynamics} ($\forall\,(i,\ell)$ s.t.\ model$=\Gamma$)
\begin{alignat}{2}
  &\alpha_{i\ell k} \geq \alpha_{i\ell,k-1} + \delta\alpha_{i\ell k}
    - \bar{\alpha}_{i\ell}\,x^m_{i\ell k}
    - \bar{\alpha}_{i\ell}\,x^r_{i\ell k}
  &&\quad\forall\,k \tag{$\Gamma$: shape accumulation} \\
  &\alpha_{i\ell k} \geq (1-\rho_{i\ell})\alpha_{i\ell,k-1}
    - \bar{\alpha}_{i\ell}(1-x^m_{i\ell k})
  &&\quad\forall\,k \tag{$\Gamma$: ARD1 repair (ARD1); ARA1 uses $\alpha^{\mathrm{last}}$} \\
  &\alpha_{i\ell k} \leq \alpha^{\mathrm{new}}_{i\ell}
    + \bar{\alpha}_{i\ell}(1-x^r_{i\ell k}),\quad
   \alpha_{i\ell k} \geq \alpha^{\mathrm{new}}_{i\ell}
    - \bar{\alpha}_{i\ell}(1-x^r_{i\ell k})
  &&\quad\forall\,k \tag{$\Gamma$: replacement} \\
  &\alpha^{\mathrm{last}}_{i\ell k} \geq \alpha_{i\ell,k-1}
    - \bar{\alpha}_{i\ell}(1-x^m_{i\ell k}),\quad
   \alpha^{\mathrm{last}}_{i\ell k} \leq \alpha_{i\ell,k-1}
    + \bar{\alpha}_{i\ell}(1-x^m_{i\ell k})
  &&\quad\forall\,k \tag{$\Gamma$: ARA1 $\alpha^{\mathrm{last}}$ update (i)} \\
  &\alpha^{\mathrm{last}}_{i\ell k} \geq \alpha^{\mathrm{last}}_{i\ell,k-1}
    - \bar{\alpha}_{i\ell}\,x^m_{i\ell k},\quad
   \alpha^{\mathrm{last}}_{i\ell k} \leq \alpha^{\mathrm{last}}_{i\ell,k-1}
    + \bar{\alpha}_{i\ell}\,x^m_{i\ell k}
  &&\quad\forall\,k \tag{$\Gamma$: ARA1 $\alpha^{\mathrm{last}}$ update (ii)} \\
  &\alpha_{i\ell k} \leq \hat{\alpha}_{i\ell}
  &&\quad\forall\,k \tag{$\Gamma$: reliability (precomputed)} \\
  &\alpha_{i\ell,2H} \leq \alpha_{i\ell,H}
  &&\tag{$\Gamma$: loop constraint}
\end{alignat}

%% --- CHANGE RF-031: rainflow dynamics block in full problem ---
\vspace{0.3em}
\noindent\textbf{Rainflow dynamics} ($\forall\,(i,\ell)$ s.t.\ model$=$RF; ARD1 only;
tail-bound and formulation row selection per the notation paragraph above)
\begin{alignat}{2}
  &\mu_{i\ell k} \geq \mu_{i\ell,k-1} + \delta\mu_{i\ell k} - \tau\,x^m_{i\ell k}
    - \tau\,x^r_{i\ell k}
  &&\quad\forall\,k \tag{RF: mean accumulation} \\
  &\mu_{i\ell k} \geq (1-\rho_{i\ell})\mu_{i\ell,k-1} - \tau(1-x^m_{i\ell k})
  &&\quad\forall\,k \tag{RF: ARD1 mean repair} \\
  &\mu_{i\ell k} \leq \mu^{\mathrm{new}}_{i\ell} + \tau(1-x^r_{i\ell k}),\quad
   \mu_{i\ell k} \geq \mu^{\mathrm{new}}_{i\ell} - \tau(1-x^r_{i\ell k})
  &&\quad\forall\,k \tag{RF: mean replacement} \\
  &v_{i\ell k} \geq v_{i\ell,k-1} + \delta v_{i\ell k}
    - V_{\max}^{\mathrm{RF}}\,x^m_{i\ell k}
    - V_{\max}^{\mathrm{RF}}\,x^r_{i\ell k}
  &&\quad\forall\,k \tag{RF: var.\ accumulation} \\
  &v_{i\ell k} \geq (1-\rho_{i\ell})^2 v_{i\ell,k-1}
    - V_{\max}^{\mathrm{RF}}(1-x^m_{i\ell k})
  &&\quad\forall\,k \tag{RF: ARD1 var.\ repair} \\
  &v_{i\ell k} \leq v^{\mathrm{new}}_{i\ell} + V_{\max}^{\mathrm{RF}}(1-x^r_{i\ell k}),\quad
   v_{i\ell k} \geq v^{\mathrm{new}}_{i\ell} - V_{\max}^{\mathrm{RF}}(1-x^r_{i\ell k})
  &&\quad\forall\,k \tag{RF: var.\ replacement} \\
  &\kappa_C^2\, v_{i\ell k} \leq (\tau - \mu_{i\ell k})^2,\quad
   \mu_{i\ell k} \leq \tau
  &&\quad\forall\,k \tag{RF: reliability (Cantelli, exact)} \\
  &\kappa_C^2\, v_{i\ell k} + 2\tau\,\mu_{i\ell k} \leq \tau^2
  &&\quad\forall\,k \tag{RF: reliability (Cantelli, LP, inner)} \\
  &(\tau - c_{i\ell} - \mu_{i\ell k})^2 \geq c_{i\ell}^2 + 2\kappa\,v_{i\ell k},\quad
   \mu_{i\ell k} \leq \tau - c_{i\ell}
  &&\quad\forall\,k \tag{RF: reliability (Bernstein, exact)} \\
  &2(\tau - c_{i\ell})\,\mu_{i\ell k} + 2\kappa\,v_{i\ell k}
    \leq \tau(\tau - 2c_{i\ell})
  &&\quad\forall\,k \tag{RF: reliability (Bernstein, LP, inner)} \\
  &\mu_{i\ell,2H} \leq \mu_{i\ell,H},\quad v_{i\ell,2H} \leq v_{i\ell,H}
  &&\tag{RF: loop constraint}
\end{alignat}

\vspace{0.3em}
\noindent\textbf{Repair cost} ($\forall\,i,\ell,k$; $\mu_{i\ell,k-1}=\alpha_{i\ell,k-1}\beta_{i\ell}$ for $\Gamma$; ARD1 form shown)
\begin{alignat}{2}
  &z_{i\ell k} \leq \rho_{i\ell}\,\mu_{i\ell,k-1}
  &&\tag{repair proxy (i)} \\
  &z_{i\ell k} \leq \rho_{i\ell}\,\tau\,x^m_{i\ell k}
  &&\tag{repair proxy (ii)} \\
  &z_{i\ell k} \geq \rho_{i\ell}\,\mu_{i\ell,k-1} - \rho_{i\ell}\,\tau\,(1-x^m_{i\ell k})
  &&\tag{repair proxy (iii)} \\
  &z_{i\ell k} \geq 0
  &&\tag{repair proxy (iv)}
\end{alignat}

\vspace{0.3em}
\noindent\textbf{Damage regularisation} (\texttt{inf\_norm}; $\mu_{i\ell k}=\alpha_{i\ell k}\beta_{i\ell}$ for $\Gamma$)
\begin{alignat}{2}
  &u \geq \textstyle\sum_{\ell=1}^{L} \mu_{i\ell k}
  &&\quad\forall\,i,k \tag{damage reg.} \\
  &u \geq 0
  &&\tag{damage reg.\ non-neg.}
\end{alignat}

\vspace{0.3em}
\noindent\textbf{Loop penalty}
\begin{alignat}{2}
  &\mathcal{L}_{\mathrm{loop}} =
    \sum_{(i,\ell):\mathrm{G}}[(\mu_{i\ell H}-\mu_{i\ell 2H})+(v_{i\ell H}-v_{i\ell 2H})]
    + \sum_{(i,\ell):\mathrm{IG}}(\mu_{i\ell H}-\mu_{i\ell 2H}) \notag \\
  &\phantom{\mathcal{L}_{\mathrm{loop}} =\;}
    + \sum_{(i,\ell):\mathrm{W}}[(\mu_{i\ell H}-\mu_{i\ell 2H})+(v_{i\ell H}-v_{i\ell 2H})]
    + \sum_{(i,\ell):\Gamma}(\alpha_{i\ell H}-\alpha_{i\ell 2H}) \notag \\
  %% --- CHANGE RF-032: rainflow term in loop penalty ---
  &\phantom{\mathcal{L}_{\mathrm{loop}} =\;}
    + \sum_{(i,\ell):\mathrm{RF}}[(\mu_{i\ell H}-\mu_{i\ell 2H})+(v_{i\ell H}-v_{i\ell 2H})]
  &&\tag{loop penalty}
\end{alignat}

\vspace{0.3em}
\noindent\textbf{Variable domains}
\begin{alignat}{2}
  &x_{ijk} \in \{0,1\},\quad x^m_{i\ell k} \in \{0,1\},\quad x^r_{i\ell k} \in \{0,1\}
  &&\tag{binary variables} \\
  &\mu_{i\ell k},\, v_{i\ell k},\, \alpha_{i\ell k},\, z_{i\ell k},\,
   \mu^{\mathrm{last}}_{i\ell k},\,
   \alpha^{\mathrm{last}}_{i\ell k},\, u \geq 0
  &&\tag{non-negativity}
\end{alignat}

\end{small}
\newpage

% ================================================================
\section{Output Specification}
% ================================================================

\subsection{Single-Horizon Output ($H$ scalar)}

The \texttt{solve()} function returns a dictionary with:

%% --- CHANGE MATH-022: revised output spec ---
\begin{longtable}{@{}llp{7cm}@{}}
\toprule
Key & Shape & Description \\
\midrule
\texttt{status} & \texttt{str} & Gurobi solver status (\texttt{"optimal"} or status code) \\
\texttt{objective} & \texttt{float} & Optimal objective value (or \texttt{null}) \\
\texttt{H} & \texttt{int} & Half-horizon used for this solve \\
\texttt{x} & $F \times (M{+}1) \times 2H$ & Assignment solution \\
\texttt{x\_m} & $F \times L \times 2H$ & Maintenance binary solution \\
\texttt{x\_r} & $F \times L \times 2H$ & Replacement binary solution \\
\texttt{mu} & $F \times L \times 2H$ & Mean degradation trajectory \\
\texttt{v} & $F \times L \times 2H$ & Variance trajectory. For components or models
  where variance is not tracked (IG, Gamma), the corresponding entries are
  \texttt{NaN}. This uniform shape matches all other per-component outputs \\
\texttt{u} & scalar & Damage regularisation variable value \\
\texttt{z} & $F \times L \times 2H$ & Degradation removed by repair in mean units \\
\texttt{model} & $F \times L$ & Model assignment (echo of input) \\
%% --- CHANGE RF-033: tail_bound echoed in output for auditability ---
\texttt{tail\_bound} & $F \times L$ list of \texttt{str} & Resolved tail bound per
  component, echoed for auditability (rainflow components only; entries for
  other components are \texttt{null}) \\
\texttt{F}, \texttt{M}, \texttt{H}, \texttt{L} & \texttt{int} & Problem dimensions \\
\bottomrule
\end{longtable}

All array outputs must be \texttt{numpy.ndarray} objects (not Gurobi objects).

\subsection{Multi-Horizon Output ($H$ interval)}

When $H = [H_{\min}, H_{\max}]$, \texttt{solve()} returns a dictionary keyed by
$H \in \{H_{\min},\ldots,H_{\max}\}$, where each value is a single-horizon result
dictionary as defined above:

\begin{lstlisting}
{
  H_min:   { "status": ..., "objective": ..., "x": ..., ... },
  H_min+1: { "status": ..., "objective": ..., "x": ..., ... },
  ...
  H_max:   { "status": ..., "objective": ..., "x": ..., ... }
}
\end{lstlisting}

Note that array shapes differ across keys since each inner dict has its own $H$.

\subsection{Solver Loop Behaviour}

\begin{lstlisting}[language=Python]
# Sequential (n_workers=1)
results = {}
hint = None
for H_k in range(H_min, H_max + 1):
    results[H_k] = solve_single(H_k, hint=hint if warm_start else None)
    if warm_start:
        hint = results[H_k]   # passed as VarHintVal; no feasibility guarantee

# Parallel (n_workers > 1); warm_start must be False
with ProcessPoolExecutor(max_workers=n_workers) as pool:
    futures = {H_k: pool.submit(solve_single, H_k) for H_k in range(H_min, H_max+1)}
    results = {H_k: fut.result() for H_k, fut in futures.items()}
\end{lstlisting}

Each parallel worker instantiates its own independent Gurobi environment.
The \texttt{verbose} and \texttt{mip\_gap} settings are shared across all solves.

%% --- CHANGE CS-005: license, error handling, trade-off guidance ---
\textbf{License requirements.}
Each parallel Gurobi solve requires a separate license token. With $N$ parallel
horizons, $N$ tokens are consumed simultaneously. For academic or single-machine
licenses (1 token), parallelism must be sequential (\texttt{n\_workers=1}).

\textbf{Error handling.}
If any horizon solve fails (infeasible, time limit, license error), the warm-start
chain is broken. The implementation shall: (1)~log the failure with the specific $H$
value, (2)~skip subsequent warm-starts that depend on the failed solve,
(3)~return partial results with a status indicating which horizons succeeded.

\textbf{Performance trade-off.}
Warm-start (sequential, reusing the previous $H$ solution as MIP start for $H+1$)
typically reduces total solve time by 30--50\% versus independent parallel solves.
Prefer warm-start unless wall-clock time is the binding constraint and multiple
license tokens are available.

% ================================================================
\section{File \& Folder Structure}
% ================================================================

%% --- CHANGE RF-034: rainflow.py added to tree; rainflow params in base.py model_params ---
\begin{lstlisting}
fleet_management/
    pyproject.toml
    README.md
    src/
        fleet_management/
            __init__.py          # Public API: solve, plot_management
            solver.py            # I/O, validation, dispatch; horizon loop
            models/
                base.py          # Abstract degradation model interface
                gaussian.py      # Gaussian component MILP builder
                inverse_gaussian.py  # IG component MILP builder
                wiener.py        # Wiener component MILP builder
                gamma.py         # Gamma component MILP builder
                rainflow.py      # Rainflow (distribution-free) component MILP builder
            maintenance/
                ard1.py          # ARD1 dynamics (shared logic)
                ara1.py          # ARA1 dynamics + big-M linearisation
            plotter.py           # Schedule visualisation (single and multi-horizon)
    input/
        data_example.yaml        # Example input (mixed models, scalar H)
        data_example_loop.yaml   # Example input (mixed models, interval H)
    results/
\end{lstlisting}

%% --- CHANGE CS-008: base.py interface with method signatures ---
\textbf{Model interface (\texttt{base.py}).}
Each degradation model module shall implement two functions with the following
signatures:

\begin{lstlisting}[language=Python]
def validate_inputs(
    F: int, H: int, M: int, L: int,
    mu_inc: np.ndarray,      # (F, M, L, 2H)
    mu_0: np.ndarray,        # (F, L)
    tau: np.ndarray,          # (F, L)
    rho: np.ndarray,          # (F, L)
    alpha_target: np.ndarray, # (F, L)
    **model_params            # model-specific (e.g., v_0, v_new, sigma, eta, beta;
                              # rainflow: v_0, v_new, b_inc, b_0, b_new, tail_bound)
) -> None:                    # raises ValueError on invalid input

def solve_fleet_management(
    F: int, H: int, M: int, L: int,
    mu_inc: np.ndarray,
    # ... same validated parameters ...
    maintenance_type: str,    # "ARD1" or "ARA1"
    mip_gap: float = 1e-4,
    time_limit: float = 3600,
    verbose: bool = False,
    warm_start: Optional[dict] = None,
) -> dict:                    # keys: status, objective, x, mu, u, z, ...
\end{lstlisting}
The \texttt{solver.py} mid-layer dispatches to the correct module based on the
\texttt{degradation\_model} input parameter. Dispatch is a simple dictionary lookup,
not class inheritance.

\subsection{API: \texttt{plot\_management}}

\begin{lstlisting}[language=Python]
plot_management(input_file_path, plot_file_path=None)
\end{lstlisting}

\begin{longtable}{@{}llp{7.5cm}@{}}
\toprule
Parameter & Type & Description \\
\midrule
\texttt{input\_file\_path} & \texttt{str} & Path to a solver output file.
  May contain a single-horizon or multi-horizon result \\
\texttt{plot\_file\_path} & \texttt{str}, optional & Output image path or \emph{prefix}.
  For multi-horizon output, one file per $H$ is produced:
  e.g. \texttt{"results/schedule"} $\to$
  \texttt{"results/schedule\_H5.png"}, \texttt{"results/schedule\_H6.png"}, \ldots
  Defaults to \texttt{"output.png"} (scalar) or \texttt{"output"} (interval) \\
\bottomrule
\end{longtable}

Supported image formats: \textbf{PNG} (\texttt{.png}), \textbf{PDF} (\texttt{.pdf}).
The format is inferred from the extension of \texttt{plot\_file\_path}; for multi-horizon
output the extension is stripped from the prefix and re-appended per file.

\subsubsection{Plot Layout}

The plot is an $F \times 2H$ grid where:
\begin{itemize}
  \item Each \textbf{row} corresponds to a train $i$.
  \item Each \textbf{column} corresponds to a time step $k$.
  \item Each cell is split into $L$ \textbf{horizontal strips}, one per component $\ell$.
  \item Each strip is coloured on a \textbf{green-to-red heatmap} based on the
        component's mean degradation $\mu_{i\ell k}$ normalised by $\tau$
        (green $= 0$, red $= \tau$). For Gamma components,
        $\mu_{i\ell k} = \alpha_{i\ell k}\cdot\beta_{i\ell}$ is used.
  \item A \textbf{thin border} around each strip indicates the degradation model:
        solid (Gaussian), dashed (IG), dotted (Wiener), dash-dot (Gamma),
        %% --- CHANGE RF-043 (panel): fifth border style for rainflow ---
        long-dash (Rainflow; e.g.\ matplotlib dash pattern \texttt{(0,\,(8,\,2))}).
\end{itemize}

Cell annotations (centred, applied to the full cell not individual strips):
\begin{itemize}
  \item \textbf{Number} $j$: mission $j$ is assigned ($x_{ijk} = 1$, $j > 0$).
  \item \textbf{Gear icon} ($\boldsymbol{\mathsf{M}}$): maintenance day
        ($x_{i0k} = 1$) with at least one repair ($x^m_{i\ell k} = 1$ for some $\ell$).
  \item \textbf{Wrench icon} ($\boldsymbol{\mathsf{R}}$): replacement scheduled
        ($x^r_{i\ell k} = 1$ for some $\ell$); may co-occur with gear icon.
  \item \textbf{``zzz''}: train is idle (no $x_{ijk} = 1$ for any $j$,
        no maintenance, no replacement).
\end{itemize}

A \textbf{vertical dashed line} separates the two half-horizons at $k = H$.
A \textbf{colourbar} is shown to the right of the grid, labelled ``Mean degradation
($\mu / \tau$)''. The plot title includes $F$, $M$, $H$, $L$, and the objective value.

% ================================================================
\section{Input File Schema}
% ================================================================

\subsection{Format Rules}

Supported formats: \textbf{YAML} (\texttt{.yaml}, \texttt{.yml}),
\textbf{JSON} (\texttt{.json}), \textbf{HDF5} (\texttt{.h5}, \texttt{.hdf5}).

Model-specific fields are \textbf{omitted entirely} when not needed (Option A):
the parser infers required fields from \texttt{model} and raises \texttt{ValueError}
if a required field is missing for any component. Fields present for non-applicable
components are ignored with a \texttt{UserWarning}.

%% --- CHANGE CS-014: HDF5 schema ---
\subsection{HDF5 Schema}

HDF5 output files shall use the following structure:
\begin{lstlisting}
/metadata/
    status          (scalar, string)
    objective       (scalar, float64)
    F, H, M, L     (scalar, int32)
    solve_time      (scalar, float64)
/solution/
    x               (int8,    shape F x (M+1) x 2H)
    mu              (float64, shape F x L x 2H)
    v               (float64, shape F x L x 2H)  # NaN if not tracked
    u               (float64, scalar)
    z               (float64, shape F x L x 2H)
/parameters/
    alpha_target    (float64, shape F x L)
    tau             (float64, shape F x L)
    rho             (float64, shape F x L)
    tail_bound      (string, scalar or shape F x L; only if rainflow components present)
    b_inc           (float64, shape F x M x L x H; rainflow + Bernstein only)
    b_0             (float64, shape F x L; rainflow + Bernstein only)
    b_new           (float64, shape F x L; rainflow + Bernstein only)
\end{lstlisting}
%% --- CHANGE RF-035: tail_bound persisted in HDF5 /parameters/ ---
Dataset names use \texttt{snake\_case}. All arrays use C-contiguous (row-major) layout.
Scalar attributes are stored as HDF5 attributes on the root group.

\subsection{Always-Present Fields}

\begin{longtable}{@{}llp{7cm}@{}}
\toprule
Field & Type / Shape & Description \\
\midrule
\texttt{F} & \texttt{int} & Number of trains \\
\texttt{H} & \texttt{int} or \texttt{[int, int]} & Half-horizon \\
\texttt{M} & \texttt{int} & Number of missions \\
\texttt{L} & \texttt{int} & Number of components per train \\
\texttt{tau} & \texttt{float} & Degradation upper bound \\
\texttt{epsilon} & \texttt{float} & Reliability threshold \\
\texttt{C\_M} & \texttt{float} & Maintenance cost coefficient \\
\texttt{C\_P} & \texttt{float} & Loop penalty coefficient \\
\texttt{C\_D} & \texttt{float} & Damage regularisation coefficient \\
\texttt{penalty\_type} & \texttt{str}, optional & \texttt{"inf\_norm"} (default)
  or \texttt{"quadratic"} \\
%% --- CHANGE RF-036: formulation token rename; tail_bound row; mu row applicability ---
\texttt{formulation} & \texttt{str}, optional & \texttt{"exact"} (default)
  or \texttt{"lp"}; \texttt{"socp"} accepted as a deprecated alias for
  \texttt{"exact"} \\
\texttt{model} & $F \times L$ list of str & Degradation model per component \\
\texttt{maintenance\_type} & $F \times L$ list of str & \texttt{"ARA1"} or \texttt{"ARD1"} \\
\texttt{rho} & $F \times L$ & Repair efficiency \\
\texttt{C\_R} & $F \times L$ & Repair cost per unit degradation removed \\
\texttt{C\_rep} & $F \times L$ & Replacement cost \\
\texttt{mu\_0} & $F \times L$ & Initial mean degradation \\
\texttt{mu\_new} & $F \times L$ & Post-replacement mean (default: 0) \\
\texttt{mu} & $F \times M \times L \times H$ & Mean increment (Gaussian, IG, Wiener, Rainflow) \\
\texttt{verbose} & \texttt{int}, optional & Gurobi verbosity (default: 1) \\
\texttt{mip\_gap} & \texttt{float}, optional & MIP gap tolerance (default: Gurobi default) \\
\texttt{n\_workers} & \texttt{int}, optional & Parallel workers (default: 1) \\
\texttt{warm\_start} & \texttt{bool}, optional & Warm start hint (default: \texttt{false}) \\
\bottomrule
\end{longtable}

\subsection{Model-Specific Fields}

Fields are required if and only if at least one component uses the corresponding
model. For the rainflow \texttt{b\_*} fields the condition is two-level, as
stated in the Required-for column: model \texttt{"rainflow"} \emph{and}
\texttt{tail\_bound="bernstein"} for at least one component.

\begin{longtable}{@{}lllp{5.5cm}@{}}
\toprule
Field & Shape & Required for & Description \\
\midrule
%% --- CHANGE RF-037: rainflow model-specific fields ---
%% --- CHANGE RF-044 (panel): tail_bound moved here from Always-Present ---
\texttt{tail\_bound} & \texttt{str} / $F {\times} L$ list & Rainflow &
  \texttt{"cantelli"} (default) or \texttt{"bernstein"}; per-component
  tail bound \\
\texttt{v\_0} & $F \times L$ & Gaussian, Wiener, Rainflow & Initial variance \\
\texttt{v\_new} & $F \times L$ & Gaussian, Wiener, Rainflow & Post-replacement variance (default: 0) \\
\texttt{v} & $F \times M \times L \times H$ & Gaussian, Rainflow & Variance increment \\
\texttt{alpha\_inc} & $F \times M \times L \times H$ & Gamma & Shape increment \\
\texttt{eta} & $F \times L$ & IG & $\mu/\lambda$ ratio \\
\texttt{sigma} & $F \times L$ & Wiener & Diffusion constant \\
\texttt{beta} & $F \times L$ & Gamma & Scale parameter \\
\texttt{b\_inc} & $F \times M \times L \times H$ & Rainflow (Bernstein) &
  Increment support upper bound \\
\texttt{b\_0} & $F \times L$ & Rainflow (Bernstein) &
  Initial-state support upper bound \\
\texttt{b\_new} & $F \times L$ & Rainflow (Bernstein) &
  Post-replacement support upper bound (default: 0) \\
%% --- CHANGE DOM-026: V_max_user in input schema ---
\texttt{v\_max\_user} & $F \times L$, optional & Gaussian, Wiener, Rainflow &
  User-defined variance upper bound \\
\bottomrule
\end{longtable}

\subsection{YAML Example (Mixed Models)}

%% --- CHANGE RF-038: rainflow component, tail_bound, b_* fields, formulation token in YAML example ---
\begin{lstlisting}
F: 4
H: 10          # or H: [8, 12] for horizon loop
M: 2
L: 2
tau: 1.0
epsilon: 0.01   # must lie in (0, 0.01]; with b_0=0.3, tau=1.0 the Bernstein
                # precondition tau > 2*kappa*b/3 = 0.921 still holds
C_M: 1.0
C_P: 2.0
C_D: 1.0
penalty_type: inf_norm   # or "quadratic" (requires formulation="exact")
formulation: exact       # or "lp"; "socp" is a deprecated alias for "exact"
n_workers: 4
warm_start: false
verbose: 0
mip_gap: 0.01

model:
  - ["gaussian", "inverse_gaussian"]   # train 1: Gaussian + IG
  - ["wiener",   "gamma"           ]   # train 2: Wiener + Gamma
  - ["gaussian", "wiener"          ]   # train 3: Gaussian + Wiener
  - ["rainflow", "inverse_gaussian"]   # train 4: Rainflow + IG

tail_bound:              # rainflow components only; scalar or F x L list
  - [null,        null]
  - [null,        null]
  - [null,        null]
  - ["bernstein", null]  # train 4: rainflow uses the Bernstein bound

maintenance_type:
  - ["ARD1", "ARD1"]    # train 1: Gaussian must be ARD1; IG must be ARD1
  - ["ARD1", "ARA1"]    # train 2: Wiener supports both; Gamma supports both
  - ["ARD1", "ARD1"]    # train 3: Gaussian must be ARD1; Wiener ARD1 here
  - ["ARD1", "ARD1"]    # train 4: Rainflow must be ARD1; IG must be ARD1

rho:
  - [0.8, 0.7]
  - [0.9, 0.6]
  - [0.75, 0.85]
  - [0.65, 0.8]

C_R:
  - [2.0, 1.5]
  - [1.8, 2.2]
  - [2.1, 1.9]
  - [1.6, 2.0]

C_rep:
  - [10.0, 8.0]
  - [9.0,  11.0]
  - [10.5, 9.5]
  - [8.5,  10.0]

mu_0:
  - [0.12, 0.10]
  - [0.15, 0.08]
  - [0.11, 0.13]
  - [0.09, 0.14]

mu_new:
  - [0.0, 0.0]
  - [0.0, 0.0]
  - [0.0, 0.0]
  - [0.0, 0.0]

# Gaussian / IG / Wiener / Rainflow: mean increment F x M x L x H
mu:
  # ... (4 x 2 x 2 x 10 tensor)

# Gaussian / Wiener / Rainflow: initial variance F x L
v_0:
  - [0.005, null]    # train 1: Gaussian needs v_0; IG does not
  - [0.006, null]    # train 2: Wiener needs v_0; Gamma does not
  - [0.004, 0.007]   # train 3: both Gaussian and Wiener need v_0
  - [0.006, null]    # train 4: Rainflow needs v_0; IG does not

# Gaussian / Rainflow: variance increment F x M x L x H
v:
  # ... (only entries for Gaussian and rainflow components)

# Wiener only: diffusion constant F x L
sigma:
  - [null, null]
  - [0.03, null]
  - [null, 0.04]
  - [null, null]

# Gamma only: shape increment F x M x L x H
alpha_inc:
  # ... (only entries for Gamma components)

# Gamma only: scale parameter F x L
beta:
  - [null, null]
  - [null, 0.5 ]
  - [null, null]
  - [null, null]

# IG only: mu/lambda ratio F x L
eta:
  - [null, 0.15]
  - [null, null]
  - [null, null]
  - [null, 0.12]

# Rainflow (Bernstein) only: increment support upper bound F x M x L x H
b_inc:
  # ... (only entries for rainflow components with tail_bound="bernstein")

# Rainflow (Bernstein) only: initial-state support upper bound F x L
b_0:
  - [null, null]
  - [null, null]
  - [null, null]
  - [0.3,  null]

# Rainflow (Bernstein) only: post-replacement support upper bound F x L (default: 0)
b_new:
  - [null, null]
  - [null, null]
  - [null, null]
  - [0.0,  null]
\end{lstlisting}

\textbf{Parser behaviour for null entries:} \texttt{null} values in model-specific
fields are ignored. The parser only reads entries $(i,\ell)$ where
\texttt{model[i][l]} matches the field's required model. Non-null values at
non-applicable positions trigger a \texttt{UserWarning} and are ignored.

% ================================================================
\section{Open Questions \& TODOs}
% ================================================================

\begin{itemize}[leftmargin=1.5em]
  %% --- CHANGE RF-045 (panel): five models ---
  \item[\checkmark] \textit{Resolved.} All five degradation models fully specified
        (Gaussian, IG, Gamma, Wiener, Rainflow) with exact and linearised forms.
  \item[\checkmark] \textit{Resolved.} Repair cost proxy: $z_{i\ell k} = \rho_{i\ell}\,\mu_{i\ell,k-1}\,x^m_{i\ell k}$
        in mean units across all models; McCormick linearisation documented.
  \item[\checkmark] \textit{Resolved.} Fleet-availability constraint removed;
        assignment constraints alone enforce correct availability.
  \item[\checkmark] \textit{Resolved.} Horizon loop: scalar or interval $H$;
        \texttt{n\_workers} parallelism; optional warm-start hinting.
  \item[\checkmark] \textit{Resolved.} Big-M values fully specified throughout.
  \item[\checkmark] \textit{Resolved.} Loop constraints: hard $+$ soft for all models;
        Wiener variance controlled via replacement (multiple replacements per horizon
        allowed; no additional feasibility check needed).
  \item[\checkmark] \textit{Resolved.} Damage regularisation: $C_S$ and per-step $u_k$
        removed; replaced by scalar $u = \max_{i,k}\sum_\ell \mu_{i\ell k}$
        (or $\sum_\ell \mu_{i\ell k}^2$) with coefficient $C_D$ and user-specified
        \texttt{penalty\_type} (\texttt{"inf\_norm"} default, \texttt{"quadratic"} option).
  \item[\checkmark] \textit{Resolved.} Plot layout fully specified: $F \times 2H$ grid,
        $L$ strips per cell, green-to-red heatmap, model border style, cell annotations,
        half-horizon separator, colourbar.
  \item[\checkmark] \textit{Resolved.} Input file schema: YAML/JSON/HDF5;
        model-specific fields omitted when not needed (Option A);
        \texttt{null} entries ignored with \texttt{UserWarning};
        mixed-model YAML example provided.
  %% --- CHANGE RF-039: problem class restated (nonconvex MIQCP, socp->exact) ---
  \item[\checkmark] \textit{Resolved.} Problem class: nonconvex MIQCP with
        \texttt{formulation="exact"} (default) when Gaussian, Wiener, or
        rainflow components are present (an IG/Gamma-only fleet remains a pure
        MILP, or a convex MIQCP if the \texttt{"quadratic"} penalty is
        selected) --- the exact reliability
        constraints of the Gaussian, Wiener, and rainflow models are
        \emph{nonconvex} quadratics (Section~\ref{sec:dyn_gaussian}), handled by
        Gurobi's nonconvex quadratic support, while the optional
        \texttt{"quadratic"} damage penalty remains a genuinely convex rotated
        SOC; pure MILP with \texttt{formulation="lp"} (all models supported).
        Gaussian/Wiener LP reliability constraint: inner approximation
        $(\Phi^{-1}(1-\varepsilon))^2 v_{i\ell k} + 2\tau\,\mu_{i\ell k} \leq \tau^2$,
        obtained by dropping $\mu_{i\ell k}^2 \geq 0$ from the exact squared
        constraint; error $= \mu_{i\ell k}^2$, small when
        $\mu_{i\ell k} \ll \tau$; rainflow LP surrogates analogous
        (Section~\ref{sec:dyn_rainflow}).
        \texttt{"quadratic"} penalty incompatible with \texttt{"lp"}.
        The legacy token \texttt{"socp"} is a deprecated alias for
        \texttt{"exact"}.
  \item[\textbf{FUTURE}] \textbf{Hierarchical failure probability bounds.}
        Currently $\varepsilon$ is a single scalar bounding the probability of failure
        of each individual component: $\Pr(D_{i\ell k} > \tau) \leq \varepsilon$
        $\forall\,i,\ell,k$. In the future it may be desirable to allow the user to
        specify the reliability requirement at a coarser level:
        \begin{itemize}
          \item \textbf{Train level:} $\Pr(\exists\,\ell : D_{i\ell k} > \tau) \leq \varepsilon_i$
                $\forall\,i,k$ --- the probability that \emph{any} component of train $i$
                fails exceeds threshold. Under component independence this becomes
                $1 - \prod_\ell (1 - \Pr(D_{i\ell k} > \tau)) \leq \varepsilon_i$,
                a nonlinear constraint coupling components within a train.
          \item \textbf{Fleet level:} $\Pr(\exists\,i,\ell : D_{i\ell k} > \tau) \leq \varepsilon$
                $\forall\,k$ --- the probability that \emph{any} component across the
                entire fleet fails. Under independence this couples all trains.
        \end{itemize}
        Both generalisations require either Boole/union-bound approximations
        (conservative but linear) or product-form reformulations (exact but nonlinear).
        The current component-level $\varepsilon$ remains a valid conservative proxy
        for both via the union bound: $\varepsilon \leq \varepsilon_{\mathrm{train}} / L$
        and $\varepsilon \leq \varepsilon_{\mathrm{fleet}} / (FL)$.
  %% --- CHANGE RF-040: rainflow open questions ---
  \item[\textbf{FUTURE}] \textbf{Schedule-dependent Bernstein support.} The
        Bernstein constant $b_{i\ell}$ is the maximum over \emph{all} admissible
        summands (initial state, replacement draws, and every mission
        increment). A per-day weighted running maximum---tracking
        $w_t\,b_t$ along the schedule actually chosen---would tighten the
        constant, but requires nonlinear (running-max) state tracking;
        deliberately not done in this revision.
  \item[\textbf{FUTURE}] \textbf{Multi-tangent PWL tightening of the rainflow LP
        surrogates.} The pointwise minimum of several tangents is a valid, less
        conservative over-estimator, but it must be encoded as a PWL
        \emph{equality} (SOS2 $\lambda$-formulation or
        \texttt{addGenConstrPWL}); a conjunction of tangent inequalities on an
        auxiliary variable is either over-conservative or unsafe (see the remark
        in Section~\ref{sec:dyn_rainflow}).
  \item[\textbf{FUTURE}] \textbf{ARA1 semantics for rainflow variance.} As for
        the Gaussian model (DOM-009), the effect of ARA1 repair on the variance
        state has no established basis; rainflow remains ARD1-only until this is
        resolved.
  \item[\textbf{FUTURE}] \textbf{Additional \texttt{tail\_bound} values.}
        Count-weighted Chernoff/Hoeffding certificates (per-mission cumulant
        generating functions, or per-mission supports only) are tighter deep in
        the tail and retain per-mission support resolution; they require a fixed
        conservative $s$ or a grid sweep, and the loop-on-moments argument of
        Section~\ref{sec:dyn_rainflow} would need re-examination (imperfect
        repair breaks the additivity of the cumulant at a fixed $s$).
  \item[\textbf{FUTURE}] \textbf{Sampling error of the empirical moments.}
        Rainflow counting yields \emph{estimates} of the per-mission mean,
        variance, and support, not exact values; a guaranteed reliability bound
        would require one-sided confidence margins on these inputs.
\end{itemize}

No blocking TODOs remain. The specification is complete pending implementation review.

\emph{Change log (v0.3, panel review).}
ARA1 maintenance removed from the Gaussian model (no literature basis for the
variance extension). Output \texttt{v} shape changed from $F \times V \times 2H$
to $F \times L \times 2H$ with NaN for non-variance components.
$\varepsilon$ bound tightened from $(0, 0.5)$ to $(0, 0.01]$.
IG FSD TODO replaced with verified Proposition (MLR proof).
Tightness remark refined to three-mechanism argument.

%% --- CHANGE RF-041: v0.4 changelog ---
\emph{Change log (v0.4, rainflow extension).}
Rainflow degradation model added: distribution-free, with per-mission empirical
mean/variance/support inputs from rainflow counting; reliability certified by a
Cantelli or Bernstein tail bound (new \texttt{tail\_bound} option,
\texttt{"cantelli"} default); exact nonconvex-quadratic and conservative LP
tangent formulations both specified; loop constraint imposed on both tracked
moments, with the argument that moment domination chains across repetitions
while the scalar bound comparison $U_{2H} \leq U_H$ does not; ARD1 only,
mirroring the Gaussian precedent. \texttt{formulation} token renamed
\texttt{"socp"} $\to$ \texttt{"exact"} (\texttt{"socp"} retained as a deprecated
alias with a \texttt{UserWarning}); the previous identification of the
Gaussian/Wiener exact reliability constraint as a rotated SOC was corrected---with
the variance affine in the decisions, the squared form is a nonconvex quadratic,
handled by Gurobi's \texttt{NonConvex=2} support (the \texttt{"quadratic"}
damage penalty remains genuinely convex SOC). Wiener loop-constraint wording
corrected per accepted review item A-02 (moment bounds, not FSD). Version
reconciliation: the v0.3 change log above corresponds to the revision previously
titled ``v0.2 (revised)'' on the title page.
A post-draft panel pass (markers RF-042--RF-045 and in-place corrections)
additionally fixed the loop-induction base case ($m \geq 2$ with
Eq.~\eqref{eq:rf_loop} as the $m=1$ case), the boundary wording of the
$\tau > 2\kappa b/3$ precondition, the a-fortiori centering step in the
Bernstein application, and the problem-class qualifier for IG/Gamma-only
fleets; extended the Bhatia--Davis consistency checks to the initial and
replacement states and the shape-validation list to $b^{\mathrm{inc}}, b^{0},
b^{\mathrm{new}}$; relocated \texttt{tail\_bound} to the model-specific
schema tables; added the fifth plot border style; and corrected the YAML
example's $\varepsilon$ to respect its own $(0, 0.01]$ bound.

\end{document}
